% Options for packages loaded elsewhere 
\PassOptionsToPackage{unicode}{hyperref}
\PassOptionsToPackage{hyphens}{url}
\documentclass[a4paper]{article}
\usepackage{authoraftertitle}
\setcounter{secnumdepth}{4}
\setcounter{tocdepth}{4}

% \paragraph should also have its title on a separate line from the bodytext
\makeatletter
\renewcommand\paragraph{%
   \@startsection{paragraph}{4}{0mm}%
      {-\baselineskip}%
      {.5\baselineskip}%
      {\normalfont\normalsize\bfseries}}
\makeatother

\usepackage[a4paper, total={170mm,257mm}, left=20mm, top=20mm,]{geometry}
\usepackage{xcolor}
\usepackage{amsmath,amssymb}
\usepackage{iftex}
\ifPDFTeX
  \usepackage[T1]{fontenc}
  \usepackage[utf8]{inputenc}
  \usepackage{textcomp} % provide euro and other symbols
\else % if luatex or xetex
  \usepackage{unicode-math} % this also loads fontspec
  \defaultfontfeatures{Scale=MatchLowercase}
  \defaultfontfeatures[\rmfamily]{Ligatures=TeX,Scale=1}
\fi
\usepackage{lmodern}
\ifPDFTeX\else
  % xetex/luatex font selection
\fi
% Use upquote if available, for straight quotes in verbatim environments
\IfFileExists{upquote.sty}{\usepackage{upquote}}{}
\IfFileExists{microtype.sty}{% use microtype if available
  \usepackage[final,nopatch=footnote]{microtype}
  \UseMicrotypeSet[protrusion]{basicmath} % disable protrusion for tt fonts
}{}
\makeatletter
\@ifundefined{KOMAClassName}{% if non-KOMA class
  \IfFileExists{parskip.sty}{%
    \usepackage{parskip}
  }{% else
    \setlength{\parindent}{0pt}
    \setlength{\parskip}{6pt plus 2pt minus 1pt}}
}{% if KOMA class
  \KOMAoptions{parskip=half}}
\makeatother
\usepackage{longtable,booktabs,array}
\usepackage{multirow}
\usepackage{calc} % for calculating minipage widths
% Correct order of tables after \subsubsection or \paragraph
\usepackage{etoolbox}
\makeatletter
\patchcmd\longtable{\par}{\if@noskipsec\mbox{}\fi\par}{}{}
\makeatother
% Allow footnotes in longtable head/foot
\IfFileExists{footnotehyper.sty}{\usepackage{footnotehyper}}{\usepackage{footnote}}
\makesavenoteenv{longtable}
\setlength{\emergencystretch}{3em} % prevent overfull lines
\usepackage{bookmark}
\IfFileExists{xurl.sty}{\usepackage{xurl}}{} % add URL line breaks if available
\urlstyle{same}

\usepackage{hyperref}
\hypersetup{
  colorlinks=true,
  pdfcreator={LaTeX via pandoc}
  pdfnewwindow=true
}

\usepackage{graphicx}
\graphicspath{ {./figures/} {./intermediate/}}

%\usepackage[import=witiko/diagrams@v2]{markdown}
%\usepackage{markdown}
%\usepackage{plantuml}
\usepackage{float}
%\usepackage{verbatim}
\newcommand{\secref}[1]{\ref{#1} \nameref{#1}}
\newcommand{\gitdescribe}{\input|"git describe --tags --always --dirty --match [[:digit:]].[[:digit:]].[[:digit:]]"}

% Avoid to have to use the percentage sign as format specifier for git commands
\makeatletter
    \edef\gitdate{\@@input|"git show -s --format=\@percentchar cs"}
\makeatother

\title{Nexxtender Charger Information}
\author{Frank HJ Cuypers}
\date{\today}

%
\begin{document}
%\maketitle 
\begin{center}
\vspace*{2cm}
\LARGE\textbf{\MyTitle} \\[0.8cm]
\large \MyAuthor \\[0.6cm]
Version \gitdescribe \\[0.4cm]
\gitdate \\[0.8cm]
\end{center}

\begin{abstract}
	This document is an effort to describe the Nexxtender ecosystem.
	Currently it focuses on the Bluetooth Low Energy interface of the Nexxtender
	chargers.
	This is the interface that the Nexxtender Nexxtmove app uses to interact
	with the charger.
\end{abstract}

\tableofcontents
\listoftables
\listoffigures

\section{Introduction}
\label{introduction}
\begin{itemize}
	
	\item
	      Section \secref{ecosystem_overview} gives an high level overview of the ecosystem.
	\item
	      Section \hyperref[ble_services_overview] details the BLE services provided by the Nexxtender
	      Charger to the Nexxtmove app.
\end{itemize}


\subsection{Acronyms and Definitions}
\label{acronyms_and_definitions}
\begin{longtable}[]{ p{4cm}   p{12cm} }
	\caption{Acronyms and Definitions}
	\label{table.1}  
	\tabularnewline
	\toprule\noalign{}
	Term & Definition \\
	\midrule\noalign{}
	\endfirsthead
	\toprule\noalign{}
	Term & Definition \\
	\midrule\noalign{}
	\endhead
	\bottomrule\noalign{}
	\endlastfoot
	\href{https://bluetoothle.wiki/start}{BLE} & Bluetooth Low Energy \\
	\href{https://en.wikipedia.org/wiki/Bluetooth_Low_Energy_beacon}{BLE beacon} & Dedicated class of BLE device \\
	\href{https://saascharge.com/glossary/}{CDR} & Charging Data Record or Charging Detail Record \\
	\href{https://www.cebeo.be/catalog/nl-be/products/powerdale-nexxtender-cluster-power-module-t2-plug-32a-voor-1-laadpunt--6486546}{CLUSTER} & Nexxtender Cluster power module\\
	CT & Nexxtender Current Transformer attached to the grid \\
	\href{https://en.wikipedia.org/wiki/CBOR}{CBOR} & Concise Binary Object
	Representation. See also \url{https://cbor.me/} for computations \\
	CRC/MODBUS & 16 bit CRC, also known as
	\href{https://en.wikipedia.org/wiki/Cyclic_redundancy_check}{CRC-16-IBM}
	or \href{https://en.wikipedia.org/wiki/Modbus}{CRC-16-MODBUS}. 
 	See also	\href{https://crccalc.com/?crc=123456789&method=CRC-16/MODBUS&datatype=ascii&outtype=hex}{CRC/MODBUS 
	computations} \\
	CRC\_8\_CCITT & 8 bit CRC also know as
	\href{https://en.wikipedia.org/wiki/Cyclic_redundancy_check}{CRC-8-CCITT
	or CRC-8/SMBUS}. See also
	\href{https://crccalc.com/?crc=263b130bff13b97472d0143635984d0bb76c7bb6f45dfcc381& method=CRC-8/SMBUS&datatype=hex&outtype=hex}{CRC/SMBUS  
	computations} \\
	\href{https://www.mydiego.lu/en}{Diego} & Diego has taken over the
	nexxtmove.me web site and the Nexxtmove apps from Powerdale \\
	EV & Electric Vehicle \\
	EVSE & Electric Vehicle Supply Equipment (the charging station) \\
	\href{https://software-dl.ti.com/lprf/sdg-latest/html/ble-stack-3.x/gatt.html}{GATT} & Generic Attributes \\
	\href{https://www.cebeo.be/catalog/nl-be/products/powerdale-laadpaal-nexxtender-home-1-4kw-22-kw-t2-cable-7m-rcm-bluetooth-6496973}{HOME}
	& Nexxtender wall mountable charger. Also called \emph{Nexxtender
	Home} \\
	\href{https://en.wikipedia.org/wiki/ISO/IEC_14443}{ISO/IEC 14443} &
	Standard for Proximity cards \\
	\href{https://punchthrough.com/lightblue/}{LightBlue}® & With
	LightBlue®, you can scan, connect to and browse any nearby BLE device \\
	\href{https://en.wikipedia.org/wiki/Endianness}{Little-endian} & A
	little-endian system stores the least significant byte of a word at the
	lowest memory address \\
	Lsn & Least significant nibble \\
	MOBILE 
	& Nexxtender Mobile charger. Also called \emph{Nexxtender Mobile}. Either a MOBILE BLACK or a MOBILE RED \\
	\href{https://www.cebeo.be/catalog/nl-be/products/powerdale-mobile-red-type-2-t2-3f-4kw-11kw-incl-rcm-adapter-to-schuko-plug-6486564}{MOBILE RED}
	& Nexxtender Mobile Red charger \\
	\href{https://www.cebeo.be/catalog/nl-be/products/powerdale-mobile-black-type-2-t2-1f-1-4kw-2-3kw-incl-rcm-6486563}{MOBILE BLACK}
	& Nexxtender Mobile Black charger \\
	Msn & Most significant nibble \\
	Nexxtender & Powerdale brand name for car chargers \\
	\href{https://play.google.com/store/apps/details?id=com.powerdale.nexxtender&hl=en}{Nexxtmove}
	& Android app to control the HOME, see
	\cite{GMNEXT} \\
	\href{https://play.google.com/store/apps/details?id=com.powerdale.homeinstaller&hl=en}{Nexxtender
	Installer} & Nexxtmove app for installers, see	\cite{NMIAPP}. 
	Seems to be a	super set of Nexxtmove \\
	\href{https://play.google.com/store/apps/details?id=no.nordicsemi.android.mcp&hl=en}{nRF
	Connect for Mobile} & With nRF Connect, you can scan, connect to and
	browse any nearby BLE device. 
	It also distinguishes beacons from other BLE devices. \\
	\href{https://www.brusselstimes.com/935740/belgian-charging-station-giant-powerdales-bankruptcy-leaves-users-stranded\%20Powerdale}{Powerdale}
	& Manufacturer of the Nexxtender products. 
	Out of business. See \cite{TA2024} (in Dutch). \\
	\href{https://en.wikipedia.org/wiki/Unix_time}{Unix Time} & Number of
	seconds since 1/1/1970 00:00 \\
	\href{https://www.wireshark.org/}{Wireshark} & The
	world\textquotesingle s most popular network protocol analyzer \\
\end{longtable}

\subsection{Disclaimer}
\label{disclaimer}

The authors of this document are not linked to Powerdale or any other company managing the Nexxtender ecosystem. 
The information in this document is based on the study of the external behavior of the product;
it is at partial, incomplete and unverified. 
No guarantees on its correctness can be given by the authors, so use it at your own risk.

\section{Ecosystem Overview}
\label{ecosystem_overview}

Figure \ref{illustration.1} gives an overview of the complete Nexxtender ecosystem.

%\begin{markdown}
%### test 
%
%``` mermaid {command = "mmdc --iconPacks '@iconify-json/logos'"}
%graph TD
%    A[Start] --> B{Ready?}
%    B -- Yes --> C[Run]
%    B -- No --> D[Wait]
%```
%\end{markdown}

\begin{figure}[H]
	\caption{Ecosystem overview}
	\label{illustration.1}
	\centering
	\includegraphics[scale=0.5]{NexxtenderOverview.pdf}
\end{figure}

The following elements play a role.

\begin{itemize}
	
	\item
	      In the \textbf{Power} domain:
	      
	      \begin{itemize}
	      	  
	      	\item
	      	      The \textbf{Power Grid} delivers the power P\textsubscript{Grid} to
	      	      the House and the \emph{Nexxtender Home} Charger. 
	      	      If the solar panels produce more power then the consumption of the home and the
	      	      \emph{Nexxtender Home} charger, 
	      	      the power P\textsubscript{Grid} from the grid is negative and power is injected in the grid.
	      	\item
	      	      The optional \textbf{Solar Panels} delivers the power  P\textsubscript{Solar} if the sun is shining.
	      	\item
	      	      The \textbf{House} consumes the power P\textsubscript{House} in
	      	      function of the appliances that are using power.
	      	\item
	      	      The \textbf{Nexxtender Home Charger} consumes the power
	      	      P\textsubscript{Charger} when the EV is charging.
	      	\item
	      	      The \textbf{EV} consumes P\textsubscript{EV} when it is charging.
	      	      P\textsubscript{EV} equals P\textsubscript{Charger} minus the
	      	      consumption of the \emph{Nexxtender Home} charger itself.
	      	\item
	      	      P\textsubscript{Grid} = P\textsubscript{House} +
	      	      P\textsubscript{Charger} - P\textsubscript{Solar}
	      	\item
	      	      The power can be 1 phase or 3 phases.
	      \end{itemize}
	\item
	      In the \textbf{Control} domain:
	      
	      \begin{itemize}
	      	  
	      	\item
	      	      The \textbf{Nexxtender Home Charger}
	      	      
	      	      \begin{itemize}
	      	      	    
	      	      	\item
	      	      	      Communicates over BLE with the Nexxtmove app.
	      	      	\item
	      	      	      Communicates over ISO/IEC 14443-3 with a Badge.
	      	      	\item
	      	      	      Measures P\textsubscript{Grid} with a CT to know how much
	      	      	      P\textsubscript{EV} to deliver in ECO mode.
	      	      \end{itemize}
	      	\item
	      	      The \textbf{Nexxtmove app}
	      	      
	      	      \begin{itemize}
	      	      	    
	      	      	\item
	      	      	      Communicates with the \emph{Nexxtender Home} charger and the
	      	      	      \href{http://www.Nexxtmove.me/}{www.Nexxtmove.me} website.
	      	      	\item
	      	      	      A \href{http://www.Nexxtmove.me/}{www.Nexxtmove.me} account is
	      	      	      required for the Nexxtmove app to work.
	      	      	\item
	      	      	      Allows to control the \emph{Nexxtender Home} charger. 
	      	      	       It is possible to configure the \emph{Nexxtender Home} charger, 
	      	      	      add badges, start/stop a charging session etc.
	      	      	\item
	      	      	      Synchronizes data, including data on charging sessions, 
	      	      	      to the \href{http://www.Nexxtmove.me/}{www.Nexxtmove.me} website.
	      	      	\item
	      	      	      An \emph{Nexxtmove Installer app} exists that is used by the
	      	      	      installer installing the \emph{Nexxtender Home} charger. 
	      	      	      It has some more functionality then the regular Nexxtmove app that can be
	      	      	      used to do the initial setup of the system. 
	      	      	      It is also the \emph{Nexxtmove Installer app} that knows how to generate the BLE
	      	      	      pairing PIN from the \emph{Nexxtender Home} charger\textquotesingle s PN and SN. 
	      	      	      The \emph{Nexxtmove Installer app} requires a special Installer account.
	      	      \end{itemize}
	      	\item
	      	      The \href{http://www.nexxtmove.me/}{\textbf{www.Nexxtmove.me}}
	      	      website
	      	      
	      	      \begin{itemize}
	      	      	    
	      	      	\item
	      	      	      Requires an account
	      	      	\item
	      	      	      Controls the Nexxtmove app which in turn controls the
	      	      	      \emph{Nexxtender Home} charger.
	      	      	\item
	      	      	      Receives all information about charging events on the
	      	      	      \emph{Nexxtender Home} charger.
	      	      	\item
	      	      	      If the \emph{Nexxtender Home} charger is provided by an employer to an employee, 
	      	      	      \href{http://www.nexttmove.me/}{www.Nexxtmove.me}
	      	      	      controls the refunding of the employee by the employer.
	      	      	\item
	      	      	      Several providers of chargers rebrand this website. 
	      	      	      For instance if your employer provided you with a \emph{Nexxtender Home}
	      	      	      charger via Q8, 
	      	      	      the Q8 electric website is actually \href{http://www.nexttmove.me/}{www.Nexxtmove.me}.
	      	      \end{itemize}
	      	\item
	      	      The \textbf{Badge}
	      	      
	      	      \begin{itemize}
	      	      	    
	      	      	\item
	      	      	      Needs to be programmed before it can be used. 
	      	      	      Programming involves the use of the Nexxtmove app and the \emph{Nexxtender Home} charger.
	      	      	\item
	      	      	      Starting a charging session on the \emph{Nexxtender Home} charger
	      	      	      with a badge is possible, 
	      	      	      even without a \href{http://www.Nexxtmove.me/}{www.Nexxtmove.me} account,
	      	      	      provided that the badge was already configured with an account.
	      	      	\item
	      	      	      The badges that I have seen are simple MIFARE Classic badges. 
	      	      	      The outdated security of MIFARE Classic is not even used; 
	      	      	      only the ISO/IEC 14443-3 Unique ID of the badge is used to determine who pays.
	      	      \end{itemize}
	      	\item
	      	      The CT measures P\textsubscript{Grid} and provides the information to
	      	      the \emph{Nexxtender Home} charger.
	      \end{itemize}
\end{itemize}

In case of the \emph{Nexxtender Mobile} charger (Red or Black) , 
there is no CT to measure P\textsubscript{Grid} and no possibility to use Badges.

\section{Product Number, Serial number, PIN and pairing}
\label{product_number_serial_number_pin_and_pairing}

The \emph{Nexxtender Home} and \emph{Mobile} have a sticker on the bottom,
mentioning its product number (PN) and serial number (SN). 
This is the sole information needed by Nexxtmove app to pair with the correct device.

The product number has the format \emph{AAAAA-RR} with

\begin{itemize}
	
	\item
	      \emph{AAAAA}: a 5 digit decimal model number. 
	      This is the same number as returned by the ``Model Number String'' from table \ref{table.9}.
	\item
	      \emph{RR}: a 2 digit hexadecimal hardware revision. 
	      This is the same number as returned by the ``Hardware Revision String'' from table \ref{table.9}.
\end{itemize}

The serial number has the format \emph{YYWW-NNNNN-UU} with

\begin{itemize}
	
	\item
	      \emph{YYWW}: a 4 digit decimal that in a lot of cases looks the year
	      and week of production, example 2303.
	\item
	      \emph{NNNNN}: a 5 digit decimal sequence number. 
	      This is the same number as returned by the ``Serial Number String'' from table \ref{table.9}.
	\item
	      \emph{UU}: a 2-digit hexadecimal number
\end{itemize}

Both the Nexxtmove and Nexxtender Installer app derive from these values
an hexadecimal serial number and a 6 digit PIN. 
The Nexxtender Charger advertises its hexadecimal serial number in the Service Data of the Advertisement packet, 
see table \ref{table.2} and \ref{table.3}, 
with UUID fd47416a-95fb-4206-88b5-b4a8045f75\textbf{c0}. 
The Nexxtender Installer app connects with the BLE device that advertises the matching hexadecimal serial number. 
The derived PIN is used to complete the pairing/bonding. 
The Nexxtmove app uses the PIN computed by the Nexxtender Installer app.

\subsection{Hexadecimal serial number computation}
\label{hexadecimal_serial_number_computation}

The Hexadecimal serial number is obtained by converting the ASCII strings YY, 
WW and NNNNN (9 decimals characters in total) to the hexadecimal numbers 
YY\textquotesingle, WW\textquotesingle{} and NNNN\textquotesingle. 
YY\textquotesingle{} and WW\textquotesingle{} are then each coded on 1 byte, 
and NNNN\textquotesingle{} is coded on 2 bytes in Big Endian. 
The result is a 4-byte serial number YY\textquotesingle WW\textquotesingle NNNN\textquotesingle. 
The result for an SN of 2304-00006-E3 is 0x17040006.

\begin{figure}[H]
	\caption{Hexadecimal serial number computation}
	\label{illustration.2}
	\centering
	\includegraphics[scale=0.5]{HexadecimalSerialNumber.pdf}
\end{figure}

The major and minor number in the Apple iBeacon advertisement packet
(see table \ref{table.3}) are made up of respectively the 2 MSB
or 2 LSB of the hexadecimal serial number.

\subsection{PIN calculation}
\label{pin_calculation}

PIN calculation is explained in the private Github project \cite{GP}.

\section{Advertisement and scanning}
\label{advertisement_and_scanning}

When not connected to a client, 
BLE devices periodically send advertisement packets containing some information on their functionality. 
The Nexxtender Charger advertises sufficient information to identify it, 
including a dedicated Service Data UUID identifying it as a Nexxtender Charger, 
with the hexadecimal serial number from section
\secref{hexadecimal_serial_number_computation} as service data.

The Nexxtender Charger advertises itself also as an Apple 
iBeacon\footnote{Android does not include beacons when scanning for BLE devices, unless the app has the location permission}.
The beacon information is coded within the 0xFF Type (Manufacturer Specific Data).

So when searching for a Nexxtender Charger with a specific SN, the BLE client can either

\begin{itemize}
	
	\item
	      scan for the BLE device advertising the specific Nexxtender Charger
	      Service Data UUID with the specific hexadecimal serial number matching
	      the device
	\item
	      scan for the BLE device advertising itself as a beacon with the
	      specific proximity UUID and the Major and Minor number matching the
	      device
\end{itemize}

When the MAC of the device is already known (using a scanner like LightBlue or nRF Connect), 
the MAC can be used to directly connected with the device.

Table \ref{table.2} shows a raw advertisement packet from a 
\emph{Nexxtender Home} as it can be seen using LightBlue.

\begin{longtable}[]{p{2cm}   p{14cm}}
	\caption{Advertisement}
	\label{table.2}
	\tabularnewline
	\toprule\noalign{}
	Properties  & Values \\
	\midrule\noalign{}
	\endfirsthead
	\toprule\noalign{}
	Properties  & Values \\
	\midrule\noalign{}
	\endhead
	\bottomrule\noalign{}
	\endlastfoot
	Adv. Packet &        
	0x0201061AFF4C000215ECB0583AD71B43D2A4FDD5367063032D \textless MMMM\textgreater \textless mmmm\textgreater C80509484F4D451521C0755F04A8B4B5880642FB956A4147FD \textless MMMM\textgreater\textless mmmm\textgreater00000000 \\
\end{longtable}

The advertisement packet is a sequence of LTV (length, Type, Value) triplets, 
called Adn in\cite{NSBLE}. 
For a complete list of AD types, see section 2.3 \emph{Common Data Types} in \cite{BLSAN}. 
For a description of the 0xFF type, see \cite{KIIB}.

Splitting the advertisement packet in its components according to these
documents results in table \ref{table.3}.

\begingroup
\renewcommand*{\arraystretch}{1.4} % Makre row separation bigger, just for this table
\begin{longtable}[]{ p{0.8cm}  p{0.8cm}  p{0.1cm}  p{4cm} p{9cm} }
	\caption{Advertisement packet data structure decoding}
	\label{table.3} 
	\tabularnewline
	\toprule\noalign{}
	Length & Type & \multicolumn{2}{l}{Value} & Description \\
	\midrule\noalign{}
	\endfirsthead
	\toprule\noalign{}
	Length & Type & \multicolumn{2}{l}{Value} & Description \\
	\midrule\noalign{}
	\endhead
	\bottomrule\noalign{}
	\endlastfoot
	0x02 & &  \multicolumn{2}{l}{} & \\ 
	& 0x01 & \multicolumn{2}{l}{} & Flags \\
	& & \multicolumn{2}{l}{0x06} & \begin{itemize}
	\item
	0x04: BR/EDR not supported
	\item
	0x02: LE General Discoverable Mode supported
	\end{itemize} \\
	0x1A & & \multicolumn{2}{l}{} & \\
	& 0xFF & \multicolumn{2}{l}{} & Manufacturer Specific data \\
	& & \multicolumn{2}{l}{0x4C00} & Company identifier. 0x004C: Apple Inc (coded in little endian) \\
	& & \multicolumn{2}{l}{0x02} & Subtype: 0x02 (Apple\textquotesingle s iBeacon) \\
	& & \multicolumn{2}{l}{0x15} & Subtype length (0x15) \\
	  &   &   & 0xECB0583AD71B43D2 A4FDD5367063032D & Proximity UUID ecb0583a-d71b-43d2-a4fd-d5367063032d                       \\
	  &   &   & 0xMMMM                              & Major. The 2 MSB from the hexadecimal serial number are used, see section 
	\secref{hexadecimal_serial_number_computation}. \\
	  &   &   & 0xmmmm                              & Minor. The 2 MSB from the hexadecimal serial number are used, see section 
	\secref{hexadecimal_serial_number_computation} \\
	  &   &   & 0xC8                                & Calibrated RSSI @ 1m                                                      \\
	0x05 & & \multicolumn{2}{l}{} & \\
	& 0x09 & \multicolumn{2}{l}{} & Complete Local Name in ASCI \\
	& & \multicolumn{2}{l}{0x484F4D45} & ``HOME'' \\
	0x15 & & \multicolumn{2}{l}{} & \\
	& 0x21 & \multicolumn{2}{l}{} & Service Data - 128-bit UUID \\
	& & \multicolumn{2}{p{4cm}}{0xC0755F04A8B4B588 0642FB956A4147FD} & Service:
	fd47416a-95fb-4206-88b5-b4a8045f75\textbf{c0} coded in little endian \\
	& & \multicolumn{2}{l}{0xMMMMmmmm} & Data. This is the serial number of the device, converted
	to hexadecimal, see section \secref{hexadecimal_serial_number_computation}. \\
	0x00 & & \multicolumn{2}{l}{} & \\
	0x00 & & \multicolumn{2}{l}{} & \\
	0x00 & & \multicolumn{2}{l}{} & \\
	0x00 & & \multicolumn{2}{l}{} & \\
\end{longtable}
\endgroup

The Advertisement package has the same structure on a \emph{Nexxtender Mobile}. 
The only difference is that the value for the \emph{Complete Local Name} is 

\begin{itemize}
	\item``Mobil'' or 0x4D6F62696C65 on a Mobile Black
	\item ``MobRED'' or 0x4D6F62524544 on a Mobile Red
\end{itemize}

Table \ref{table.4} shows the UUID\textquotesingle s that are used in the advertisement packet.

\begingroup
\renewcommand*{\arraystretch}{1.4} % Makre row separation bigger, just for this table
\begin{longtable}[]{p{6.5cm} p{9cm}}
	\caption{UUID\textquotesingle s used in the advertisement package}
	\label{table.4}
	\tabularnewline
	\toprule\noalign{}
	UUID                                          & Name                  \\
	\midrule\noalign{}
	\endfirsthead
	\toprule\noalign{}
	UUID                                          & Name                  \\
	\midrule\noalign{}
	\endhead
	\bottomrule\noalign{}
	\endlastfoot
	fd47416a-95fb-4206-88b5-b4a8045f75\textbf{c0} & Service Data in table 
	\ref{table.2} and table \ref{table.3}. 
	Note that this UUID is the same as the one in table \ref{table.74},
	except for the last byte that is 0xC0. \\
	ecb0583a-d71b-43d2-a4fd-d5367063032d          & Proximity UUID        \\
\end{longtable}
\endgroup

\section{BLE services overview}
\label{ble_services_overview}

This section describes the BLE GATT Services and Characteristics
supported by the Nexxtender Charger.

\begin{itemize}
	
	\item
	      Section \secref{general_information} provides some general information.
	\item
	      Section \secref{services} gives an overview of the supported services.
	\item
	      Section \secref{characteristics} gives an overview of the supported characteristics.
	\item
	      Section \secref{generic_access_characteristics} discusses the characteristics supported by the
	      \emph{Generic Access} service.
	\item
	      Section
	      \secref{generic_attribute_characteristics} discusses the characteristics supported by
	      the \emph{Generic Attributes} service.
	\item
	      Section
	      \secref{device_information_characteristics} discusses the characteristics supported
	      by the \emph{Device Information} service.
	\item
	      Section
	      \secref{generic_characteristics} discusses the characteristics supported by the
	      \emph{Generic} service.
	\item
	      Section \secref{cdr_characteristics} discusses the characteristics supported by the
	      \emph{CDR} service.
	\item
	      Section
	      \secref{ccdt_characteristics} discusses the characteristics supported by the
	      \emph{CCDT} service.
	\item
	      Section
	      \secref{firmware_characteristics} discusses the characteristics supported by the
	      \emph{Firmware} service.
	\item
	      Section
	      \secref{charging_characteristics} discusses the characteristics supported by the
	      \emph{Charging} service.
\end{itemize}

\subsection{General information}
\label{general_information}

BLE GATT codes fields always in little-endian. 
So all integer value types wider than uint8 (uint16, uint32, \ldots) always store the least
significant byte (LSB) at the lowest address, 
then the next least significant byte, 
until the most significant byte (MSB) which is stored at the highest address. 
The only exception is the CBOR format which uses big-endian.

The Bluetooth Generic Attribute Profile (GATT) defines how data is
organized and exchanged between Bluetooth Low Energy (BLE) devices. 
It uses the Attribute Protocol (ATT) to structure data into attributes,
characteristics, and services. 
GATT profiles are standardized sets of services and characteristics that define specific use cases, 
like heart rate monitoring or HID (Human Interface Devices).

A profile consists of services that each consist of one or more characteristics. 
Characteristics can contain descriptors that give more information on the characteristic, 
like the Client Characteristic Configuration Descriptor (CCCD) 0x2902 that defines if the
characteritics supports read/write and/or notification , etc...)

The Nexxtender charger devices seem to define a proprietary ``Nexxtender Profile'' 
that uses the services and characteristics from the following diagram. 
Descriptors are not shown in this diagram. 
The greyed-out services and characteristics are standardized by the Bluetooth Special
Interest Group (SIG) and are included by the ``Nexxtender Profile''
as shown in figure \ref{illustration.25}.

\begin{figure}[H]
	\caption{Proprietary Nexxtender profile}
	\label{illustration.25}
	\centering
	\includegraphics[scale=0.5]{NexxtenderProfile.png}
\end{figure}

These services and characteristics are explained in the next sections.

The Nexxtender chargers that I have seen support a variant of this ``Nexxtender Profile''.

The \emph{Nexxtender Home} charger (``HOME'' or ``HOME2\_'') groups all Nexxtender
specific characteristics under 1 \textbf{Generic} service as shown in figure \ref{illustration.26}.

\begin{figure}[H]
	\caption{Proprietary Nexxtender Home profile}
	\label{illustration.26}
	\centering
	\includegraphics[scale=0.5]{NexxtenderHomeProfile.png}
\end{figure}

The \emph{Nexxtender Mobile} chargers (Red and Black) use specific
Nexxtender services as shown in figure \ref{illustration.27}.

\begin{figure}[H]
	\caption{Proprietary Nexxtender Mobile profile}
	\label{illustration.27}
	\centering
	\includegraphics[scale=0.5]{NexxtenderMobileProfile.png}
\end{figure}

BLE GATT services and characteristics are addressed using UUIDs. 
The server builds a GATT table with all UUIDs and assigns each entry a 16-bit handle. 
The first time that a GATT client connects to a GATT server, 
the client synchronizes with the GATT table so that it has a copy. 
From then on all GATT operations (read, write, notify, \ldots)
addresses a characteristic using its handle and not its UUID. 
Clients can cache this GATT table so that after a reconnect, 
no table synchronization is needed anymore. 
In the BLE protocol between client and server, only handles are visible. 
Protocol analyzers like Wireshark are normally able to map handles to UUIDs and show these in the GUI, 
but only if the GATT table synchronization between client and server was also part of the log; 
in most cases it is not because it was cached. 
The GATT table is static, unless the server firmware changes. 
In the tables below that define the characteristics of the Nexxtender Charger, 
a column \emph{Handle} is added with the handle values as observed on a firmware revision 2.53.2.

\subsection{Services}
\label{services}

Tables \ref{table.5} and \ref{table.74} show the GATT services that are used.

\begin{longtable}[]{@{}ll@{}}
	\caption{Bluetooth SIG predefined GATT services used}
	\label{table.5}
	\tabularnewline
	\toprule\noalign{}
	UUID                                          & Name               \\
	\midrule\noalign{}
	\endfirsthead
	\toprule\noalign{}
	UUID                                          & Name               \\
	\midrule\noalign{}
	\endhead
	\bottomrule\noalign{}
	\endlastfoot
	0000\textbf{1800}-0000-1000-8000-00805f9b34fb & Generic Access     \\
	0000\textbf{1801}-0000-1000-8000-00805f9b34fb & Generic Attribute  \\
	0000\textbf{180a}-0000-1000-8000-00805f9b34fb & Device Information \\
\end{longtable}

\begingroup
\renewcommand*{\arraystretch}{1.4} % Makre row separation bigger, just for this table
\begin{longtable}[]{@{}ll@{}}
	\caption{Nexxtender specific GATT services used}
	\label{table.74}
	\tabularnewline
	\toprule\noalign{}
	UUID                                                           & Name              \\
	\midrule\noalign{}
	\endfirsthead
	\toprule\noalign{}
	UUID                                                           & Name              \\
	\midrule\noalign{}
	\endhead
	\bottomrule\noalign{}
	\endlastfoot
	fd47416a-95fb-4206-88b5-b4a8045f75\textbf{c0}                  &                   
	\textbf{Advertisement} \\
	\multirow{2}{*}{fd47416a-95fb-4206-88b5-b4a8045f75\textbf{c1}} &                   
	\textbf{Generic} \\
	                                                               & \textbf{CDR}      \\
	fd47416a-95fb-4206-88b5-b4a8045f75\textbf{c5}                  & \textbf{CCDT}     \\
	fd47416a-95fb-4206-88b5-b4a8045f75\textbf{c9}                  & \textbf{Firmware} \\
	fd47416a-95fb-4206-88b5-b4a8045f75\textbf{ce}                  & \textbf{Charging} \\
\end{longtable}
\endgroup
The following sections specify for each service which characteristics are defined.

\subsection{Characteristics}
\label{characteristics}

The following operations can be performed by a client on a server characteristic:

\begin{itemize}
	
	\item
	      \textbf{Read} (R): the server returns the value of the characteristic.
	\item
	      \textbf{Write} (W): the server changes the value of the characteristic
	      to the one given by the client and returns an error/success
	      code\textquotesingle{}
	\item
	      \textbf{Write Without Response} (WWR): like Write, but without response.
	\item
	      \textbf{Notify} (R): server sends the new characteristic value to the
	      client as soon as it changes. 
	      Also supports \textbf{Indicate} which requires an acknowledgment by the 
	      client of the reception of the indication. 
	      Notifications must be enabled explicitly per characteristic. 
	      That is done by writing the \textbf{Client characteristic configuration descriptor (CCCD)} 
	      for each characteristic for which the client wants to receive a notification.
	      That is normally hidden/performed by the BLE library used on the client.
\end{itemize}

Not all operations can be performed on each characteristic. 
Some operations require a pairing with the PIN.
Other operations work without pairing.
The tables in the following section indicate which are available:
\begin{itemize}
	\item An empty cell means the operation is not available.
	\item A \emph{P} means the operation is available and requires a PIN verification.
	\item A \emph{N} means the operation is available without PIN verification.
	\item A \emph{P?} means the operation is available and probably requires a PIN verification (not tested).
\end{itemize}

\subsection{Generic Access characteristics}
\label{generic_access_characteristics}

Table \ref{table.6} shows the following Characteristics of the \emph{Generic Access} service that are used.

\begin{longtable}[]{p{1cm} p{5.7cm} p{3.7cm} p{0.4cm} p{0.5cm} p{0.9cm} p{0.4cm}  p{0.9cm}}
	\caption{GATT Generic Access service characteristics used}
	\label{table.6}
	\tabularnewline
	\toprule\noalign{}
	Handle & UUID & Name & R & W & WWR & N & Coding \\
	\midrule\noalign{}
	\endfirsthead
	\toprule\noalign{}
	Handle & UUID & Name & R & W & WWR & N & Coding \\
	\midrule\noalign{}
	\endhead
	\bottomrule\noalign{}
	\endlastfoot
	0x0007 & \small 0000\textbf{2a00-}0000-1000-8000-00805f9b34fb & \textbf{Device
	Name} & N & P? & P? & & ASCII \\
	0x0009 & \small 0000\textbf{2a01-}0000-1000-8000-00805f9b34fb &
	\textbf{Appearance} & N & P? & P? & & HEX \\
	0x000c & \small 0000\textbf{2a04-}0000-1000-8000-00805f9b34fb &
	\textbf{Peripheral Preferred Connection} & N? & & & & HEX \\
\end{longtable}

Note that the app must be \href{https://developer.android.com/reference/android/Manifest.permission.html#BLUETOOTH_PRIVILEGED}{ BLUETOOTH\_PRIVILEGED}.

\subsubsection{Device Name characteristic}
\label{device_name_characteristic}

The \emph{Nexxtender Home} returns the ASCII string ``HOME2'' or  ``HOME2\_''.
The \emph{Nexxtender Mobile Black} returns the ASCII string ``MOBILE\_''.
The \emph{Nexxtender Mobile Red} returns the ASCII string ``TBD''.

\subsubsection{Appearance characteristic}
\label{appearance_characteristic}

The Nexxtender Charger returns the HEX array ``00 00''.

\subsubsection{Peripheral Preferred Connection Parameters characteristic}
\label{peripheral_preferred_connection_parameters_characteristic}

The Nexxtender Charger returns the HEX array ``FF FF FF FF 00 00 FF FF''.

\subsection{Generic Attribute characteristics}
\label{generic_attribute_characteristics}

Table \ref{table.7} shows the Characteristics of the \emph{Generic Attribute} service that are used.

\begin{longtable}[]{@{}lllllll@{}}
	\caption{GATT Generic Attribute service characteristics used}
	\label{table.7}
	\tabularnewline
	\toprule\noalign{}
	UUID & Name & R & W & WWR & N & Coding \\
	\midrule\noalign{}
	\endfirsthead
	\toprule\noalign{}
	UUID & Name & R & W & WWR & N & Coding \\
	\midrule\noalign{}
	\endhead
	\bottomrule\noalign{}
	\endlastfoot
	0000\textbf{2a05-}0000-1000-8000-00805f9b34fb & \textbf{Service Changed}
	& & & & N & \\
\end{longtable}

\subsubsection{Service Changed characteristic}
\label{service_changed_characteristic}

Apparently not readable on the Nexxtender Charger, only subscribable.

\subsection{Device Information characteristics}
\label{device_information_characteristics}

Table \ref{table.8} shows the Characteristics of the \emph{Device Information} service that are used.

\begin{longtable}[]{p{1cm} p{5.7cm} p{4.5cm} p{0.3cm} p{0.3cm} p{0.7cm} p{0.3cm}  p{0.8cm}}
	\caption{GATT Device Information service characteristics used}
	\label{table.8}
	\tabularnewline
	\toprule\noalign{}
	Handle & UUID & Name & R & W & WWR & N & Coding \\
	\midrule\noalign{}
	\endfirsthead
	\toprule\noalign{}
	Handle & UUID & Name & R & W & WWR & N & Coding \\
	\midrule\noalign{}
	\endhead
	\bottomrule\noalign{}
	\endlastfoot
	0x000e & \small 0000\textbf{2a24-}0000-1000-8000-00805f9b34fb & \textbf{Model
	Number String} & N & P & & & ASCII \\
	0x0010 & \small 0000\textbf{2a25-}0000-1000-8000-00805f9b34fb & \textbf{Serial
	Number String} & N & P & & & ASCII \\
	0x0012 & \small 0000\textbf{2a26-}0000-1000-8000-00805f9b34fb &
	\textbf{Firmware Revision String} & N & P & & & ASCII \\
	0x0014 & \small 0000\textbf{2a27-}0000-1000-8000-00805f9b34fb &
	\textbf{Hardware Revision String} & N & P & & & ASCII \\
\end{longtable}

Each string is 0x00-terminated and maximum 16 bytes in length. 
An example of the returned values is in table \ref{table.9}.

\begin{longtable}[]{@{}lll@{}}
	\caption{GATT Device Information service characteristics Example values}
	\label{table.9}
	\tabularnewline
	\toprule\noalign{}
	Name                 & HEX                                             & ASCII \\
	\midrule\noalign{}
	\endfirsthead
	\toprule\noalign{}
	Name                 & HEX                                             & ASCII \\
	\midrule\noalign{}
	\endhead
	\bottomrule\noalign{}
	\endlastfoot
	Model Number String  & 36 30 32 31 31 00 00 00 00 00 00 00 00 00 00 00 &       
	60211 \\
	Serial Number String & 30 30 30 31 36 00 00 00 00 00 00 00 00 00 00 00 &       
	00016 \\
	Firmware Revision String & 32 2E 35 33 2E 32 00 00 00 00 00 00 00 00 00
	00 & 2.53.2 \\
	Hardware Revision String & 41 32 00 00 00 00 00 00 00 00 00 00 00 00 00
	00 & A2 \\
\end{longtable}

\subsection{Generic characteristics}
\label{generic_characteristics}

The Generic service is only available on the \emph{Nexxtender Home}, 
not on the \emph{Nexxtender Mobile}. 
So none of its Characteristics are available on the \emph{Nexxtender Mobile}.

Table \ref{table.10} shows the Characteristics of the \emph{Generic} service that are used by
the \emph{Nexxtender Home}.

\begin{longtable}[]{@{}llllllll@{}}
	\caption{GATT Generic service characteristics use}
	\label{table.10}
	\tabularnewline
	\toprule\noalign{}
	Handle & UUID & Name & R & W & WWR & N & Coding \\
	\midrule\noalign{}
	\endfirsthead
	\toprule\noalign{}
	Handle & UUID & Name & R & W & WWR & N & Coding \\
	\midrule\noalign{}
	\endhead
	\bottomrule\noalign{}
	\endlastfoot
	0x003f & fd47416a-95fb-4206-88b5-b4a8045f75\textbf{dd} & \textbf{Generic
	Command} & & P & & & \\
	0x0041 & fd47416a-95fb-4206-88b5-b4a8045f75\textbf{de} & \textbf{Generic
	Status} & P & & & N & \\
	0x0044 & fd47416a-95fb-4206-88b5-b4a8045f75\textbf{df} & \textbf{Generic
	Data} & N & & & P & \\
\end{longtable}

\textbf{Generic Command}, \textbf{Generic Status} and \textbf{Generic Data} 
work together to perform an operation. 
An operation can trigger some action, or reads or writes configuration data. 
Each operation is identified with a 2 bytes \emph{Operation}. 
Each \emph{Operation} can be performed using these 3 registers, 
as is explained further and starts by writing an \emph{Operation} to \textbf{Generic Command}. 
The scenario is each time similar.

Figure \ref{illustration.31} shows the scenario for performing an \emph{Operation} 
that does not read/write a value.

\begin{figure}[H]
	\caption{Operation without read or write data}
	\label{illustration.31}
	\centering
	\includegraphics[scale=0.5]{OperationNoReadNoWrite.pdf}
\end{figure}

\begin{enumerate}
	
	\item
	      Client writes \textbf{Generic Command}{[}Operation{]}
	\item
	      Client is notified of a \textbf{Generic Status} change, containing the new status.
\end{enumerate}

Figure \ref{illustration.8} shows the scenario for performing an \emph{Operation} that writes a value.

\begin{figure}[H]
	\caption{Operation writing data}
	\label{illustration.8}
	\centering
	\includegraphics[scale=0.5]{OperationNoReadWrite.pdf}
\end{figure}

\begin{enumerate}
	
	\item
	      Client writes \textbf{Generic Command}{[}Operation{]}
	\item
	      Client is notified of a \textbf{Generic Status} change, containing the new status, 
	      indicating if the target is ready to receive the data to be written.
	\item
	      The client writes a \textbf{Generic Data} with the actual data that
	      must be written for \emph{Operation}.
	\item
	      Client is notified of a \textbf{Generic Status} change, containing the
	      new status, indicating if the data was successfully written.
\end{enumerate}

Figure \ref{illustration.9} shows the scenario for performing an \emph{Operation} that reads a value.

\begin{figure}[H]
	\caption{Operation reading data}
	\label{illustration.9}
	\centering
	\includegraphics[scale=0.5]{OperationReadNoWrite.pdf}
\end{figure}

\begin{enumerate}
	
	\item
	      Client writes \textbf{Generic Command}{[}Operation{]}
	\item
	      Client is notified of a \textbf{Generic Status} change, containing the
	      new status, indicating if the target is ready to return the requested
	      data.
	\item
	      The client reads \textbf{Generic Data} and gets the requested data.
\end{enumerate}

The two bytes \emph{Operation} is be made up of:

\begin{itemize}
	
	\item
	      MSB: an \emph{operation type}. 
	      An \emph{operation type} groups operations that belong to the same functionality. 
	      Supported operation types are shown in table \ref{table.11} and detailed in the indicated sections.
\end{itemize}

\begin{longtable}[]{@{}ll@{}}
	\caption{GATT Generic characteristics service operation types}
	\label{table.11}
	\tabularnewline
	\toprule\noalign{}
	Operation Type & Functionality              \\
	\midrule\noalign{}
	\endfirsthead
	\toprule\noalign{}
	Operation Type & Functionality              \\
	\midrule\noalign{}
	\endhead
	\bottomrule\noalign{}
	\endlastfoot
	0x00           & \secref{loader_operations} \\
	0x10           & \secref{event_operations}  \\
	0x20           & \secref{metric_operations} \\
	0x30           & \secref{badge_operations}  \\
	0x40           & \secref{time_operation}    \\
	0x50           & \secref{config_operations} \\
\end{longtable}

\begin{itemize}
	
	\item
	      LSB: an \emph{operation id} within the \emph{operation type}.
\end{itemize}

The two byte \emph{status} returned by \textbf{Generic Status} is also coded on 2 bytes:

\begin{itemize}
	
	\item
	      MSB: the same \emph{operation type} as for \emph{Operation.}
	\item
	      LSB: an \emph{error id} within the \emph{status.}
\end{itemize}

\subsubsection{Loader operations}
\label{loader_operations}

Loader operations can start and stop the charger. 
The procedure is shown in figure \ref{illustration.16}.

\begin{figure}[H]
	\caption{Loader operation}
	\label{illustration.16}
	\centering
	\includegraphics[scale=0.5]{LoaderOperation.pdf}
\end{figure}

\begin{enumerate}
	
	\item
	      Client writes \textbf{Generic Command}{[}Loader Operation{]}
	\item
	      Client is notified of a \textbf{Generic Status} change, containing the
	      new loader status.
\end{enumerate}

\paragraph{Loader operation ids}
\label{loader_operation_ids}

Table \ref{table.12} lists the \emph{Loader Operation} id\textquotesingle s.

\begin{longtable}[]{@{}ll@{}}
	\caption{Loader Operation ids}
	\label{table.12}
	\tabularnewline
	\toprule\noalign{}
	Operation Id & Description                                                 \\
	\midrule\noalign{}
	\endfirsthead
	\toprule\noalign{}
	Operation Id & Description                                                 \\
	\midrule\noalign{}
	\endhead
	\bottomrule\noalign{}
	\endlastfoot
	0x0001       & Start charging with the default settings, as set by mode in 
	table \ref{table.31} \\
	0x0002       & Start charging Max                                          \\
	0x0003       & Start charging Auto\footnote{What is this?}                 \\
	0x0004       & Start charging Eco                                          \\
	0x0006       & Stop charging                                               \\
\end{longtable}

These operations have no additional configuration data that is read or
written.

\paragraph{Loader status}
\label{loader_status}

Possible \emph{Loader Status} values are shown in table \ref{table.13}.

\begin{longtable}[]{@{}ll@{}}
	\caption{Loader Status codes}
	\label{table.13}
	\tabularnewline
	\toprule\noalign{}
	Status Code & Description          \\
	\midrule\noalign{}
	\endfirsthead
	\toprule\noalign{}
	Status Code & Description          \\
	\midrule\noalign{}
	\endhead
	\bottomrule\noalign{}
	\endlastfoot
	0x0001      & UNLOCKED             \\
	0x0002      & UNLOCKED\_FORCE\_MAX \\
	0x0003      & UNLOCKED\_FORCE\_ECO \\
	Other       & UNKNOWN              \\
\end{longtable}

\subsubsection{Event operations}
\label{event_operations}

This deals with BLE events different from charging events from section \secref{ccdt_characteristics}.

\paragraph{Event operation ids}
\label{event_operation_ids}

Table \ref{table.14} lists the \emph{Event Operation} id\textquotesingle s.

\begin{longtable}[]{@{}ll@{}}
	\caption{Event operation ids}
	\label{table.14}
	\tabularnewline
	\toprule\noalign{}
	Operation Id & Description    \\
	\midrule\noalign{}
	\endfirsthead
	\toprule\noalign{}
	Operation Id & Description    \\
	\midrule\noalign{}
	\endhead
	\bottomrule\noalign{}
	\endlastfoot
	0x1001       & NEXT           \\
	0x1002       & UPDATE\_STATUS \\
\end{longtable}

These operations have no additional configuration data that is read or written.

\paragraph{Event status}
\label{event_status}

Possible \emph{Event Status} values are shown in table \ref{table.15}.

\begin{longtable}[]{@{}ll@{}}
	\caption{Event status codes}
	\label{table.15}
	\tabularnewline
	\toprule\noalign{}
	Status Code   & Description                                         \\
	\midrule\noalign{}
	\endfirsthead
	\toprule\noalign{}
	Status Code   & Description                                         \\
	\midrule\noalign{}
	\endhead
	\bottomrule\noalign{}
	\endlastfoot
	0x10\emph{nn} & \emph{nn}: the number of \emph{remaining events} to 
	read \\
\end{longtable}

\paragraph{Event list operation}
\label{event_list_operation}

\textbf{Generic Command}, \textbf{Generic Status} and \textbf{Generic Data} 
work together to load event records from the device. 
The scenario is as shown in figure \ref{illustration.20}.

\begin{figure}[H]
	\caption{Event scenario}
	\label{illustration.20}
	\centering
	\includegraphics[scale=0.5]{Events.pdf}
\end{figure}

\begin{enumerate}
	
	\item
	      Client enables \textbf{Generic Status} notifications.
	\item
	      Client writes the \textbf{Generic Command}{[}Event:UPDATE\_STATUS(0x1002){]} 
	      for the first record.
	\item
	      Client is notified of a \textbf{Generic Status} change, containing the new status, 
	      ndicating 0x10\emph{nn-\/-}, with \emph{nn} the number of
	      \textbf{Generic Data} records available. 
	      If \emph{nn} \textgreater{} 0, the scenario continues at point 4. 
	     Otherwise it is complete.
	\item
	      Client reads the \textbf{Generic Data} obtaining the data of the record
	\item
	      Client writes the \textbf{Generic Command}{[}Event:NEXT(0x1001){]} for the next record. 
	      Note that probably this also indicates that the record is successfully stored by the Client, 
	      indicating that the Server can delete the record. 
	      So if you still need to sync with Nexxtmove.me, 
	      you should not perform above scenario with another app.
	\item
	      Client is notified of a \textbf{Generic Status} change, containing the new status, 
	      indicating 0x10\emph{nn-\/-}, with \emph{nn} the number of
	      \textbf{Generic Data} records available. 
          If \emph{nn} \textgreater{} 0, the scenario continues at point 4. 
	     Otherwise it is complete.
	\item
	      Client disables \textbf{Generic Status} notifications.
\end{enumerate}

\footnote{Can this be confirmed?}

The number of Event records that can be stored by the \emph{Nexxtender Home} is limited. 
The exact limit is not known. 
If the limit is reached, it is expected that older records are overwritten.

\paragraph{Generic Data format}
\label{generic_data_format}

\textbf{Generic Data} returns 20 bytes related to an event. 
Events can happen anytime, 
not only during the start/stop time frame of a \secref{cdr_characteristics}.
The 20 bytes are structured as shown in tables \ref{table.61} and \ref{table.62}.

\begin{longtable}{|l|l|l|l|l|l|l|l|l|l|l|l|l|l|l|l|l|l|l|l|}
	\caption{Generic Data data format overview}
	\label{table.61}\\  %<---- Problem
	\hline
	0 & 1 & 2 & 3 & 4 & 5 & 6 & 7 & 8 & 9 & 10 & 11 & 12 & 13 & 14 & 15 & 16 
	& 17 & 18 & 19 \\
	\hline
	\multicolumn{4}{|l|}{%
	eventTime} & unk1 & unk2 & unk3 & \multicolumn{11}{l|}{%
	data} & \multicolumn{2}{l|}{%
	crc16} \\
	\hline
\end{longtable}

\hfill\break

\begin{longtable}[]{llllp{10cm}}
	\caption{Generic Data data format detail}
	\label{table.62}
	\tabularnewline
	\toprule\noalign{}
	Offset & Length & Name      & Type       & Description                                \\
	\midrule\noalign{}
	\endfirsthead
	\toprule\noalign{}
	Offset & Length & Name      & Type       & Description                                \\
	\midrule\noalign{}
	\endhead
	\bottomrule\noalign{}
	\endlastfoot
	0      & 4      & eventTime & uint32     & Time of the event                          \\
	4      & 1      & unk1      & byte       & Unknown. Somehow codes the event type?     \\
	5      & 1      & unk2      & byte       & Unknown. Somehow codes the event type?     \\
	6      & 1      & unk3      & byte       & Unknown. Somehow codes the event type?     \\
	7      & 11     & data      & byte{[}{]} & Data that belongs to the event. The format 
	depends on the event type. \\
	18     & 2      & crc16     & uint16     & CRC-16/MODBUS over previous 18 bytes       \\
\end{longtable}

The \emph{eventType} coding is still unclear. 
By mapping the Events transferred over BLE GATT to the events shown by the Nexxtmove app or
Nexxtmove.me website, the following partial table \ref{table.63} is constructed.

\begin{longtable}[]{llllp{2.7cm}p{5.8cm}}
	\caption{eventType coding}
	\label{table.63}
	\tabularnewline
	\toprule\noalign{}
	unk1 & unk2 & unk3 & data                                   & Name                                   & Description      \\
	\midrule\noalign{}
	\endfirsthead
	\toprule\noalign{}
	unk1 & unk2 & unk3 & data                                   & Name                                   & Description      \\
	\midrule\noalign{}
	\endhead
	\bottomrule\noalign{}
	\endlastfoot
	0x02 & 0x01 & 0x02 & \footnotesize 0x0100000000000000000000 & \footnotesize PLUGIN                   &                  \\
	0x02 & 0x01 & 0x00 & \footnotesize 0x0100000000000000000000 & \footnotesize PLUGIN                   & What is the      
	difference with the other PLUGIN row? \\
	0x02 & 0x01 & 0x01 & \footnotesize 0x0000000000000000000000 & \footnotesize UNPLUGGED                &                  \\
	0x01 & 0x01 & 0x02 & \footnotesize 0x070497db82126880000000 & \footnotesize AUTHORIZED\_\ ACCESS     &                  
	Authorized access with a badge with 7-byte UID 0x0497db82126880 \\
	0x02 & 0x01 & 0x01 & \footnotesize 0x0200000000000000000000 & \footnotesize CHARGING                 &                  \\
	0x01 & 0x01 & 0x40 & \footnotesize 0x0000000000000000000000 & \footnotesize ECO?                     & During charging, 
	switched from MAX to ECO charging using the Nexxtmove app. \\
	\multirow{5}{*}{0x06} & \multirow{5}{*}{0x01} & \multirow{5}{*}{0x08} &
	\footnotesize 0x0900020000000000000000 & \multirow{5}{*}{\footnotesize FAULT (Error)} &
	\multirow{5}{5cm}{An error event. Details are coded in the data, but syntax unknown.}\\
	& & & \footnotesize 0x0900010000000000000000 \\
	& & & \footnotesize 0x0100030000000000000000 \\
	& & & \footnotesize 0x0100050000000000000000 \\
	& & & \footnotesize 0x0900040000000000000000 \\
	0x04 & 0x01 & 0x03 & \footnotesize 0x060000000a000000000000 & \footnotesize CONFIGURATION\_\ UPDATED &                  
	The
	\secref{config_1_0_format} field \emph{safe} was changed from 6 to 10. 
	The data first codes the old value 6 on 32 bits in little endian. 
	Then it codes the new value in the same way. \\
	0x04 & 0x01 & 0x03 & \footnotesize 0x0a00000006000000000000 & \footnotesize CONFIGURATION\_\ UPDATED &                  
	\secref{config_1_0_format} field \emph{safe} was changed from 10 to 6. 
	The data first codes the old value 10 on 32 bits in little endian. 
	Then it codes the new value in the same way. \\
\end{longtable}

\subsubsection{Metric operations}
\label{metric_operations}

This deals with Metrics\footnote{Which metrics?}.

\paragraph{Metrics operation ids}
\label{metrics_operation_ids}

Table \ref{table.16} lists the \emph{Metrics Operation} id\textquotesingle s.

\begin{longtable}[]{@{}ll@{}}
	\caption{Metrics Operation ids}
	\label{table.16}
	\tabularnewline
	\toprule\noalign{}
	Operation Id & Description    \\
	\midrule\noalign{}
	\endfirsthead
	\toprule\noalign{}
	Operation Id & Description    \\
	\midrule\noalign{}
	\endhead
	\bottomrule\noalign{}
	\endlastfoot
	0x2001       & NEXT           \\
	0x2002       & UPDATE\_STATUS \\
\end{longtable}

These operations have no additional configuration data that is read or written.

\paragraph{Metrics status}
\label{metrics_status}

Possible Metrics \emph{status} values are shown in table \ref{table.17}.

\begin{longtable}[]{@{}ll@{}}
	\caption{Metrics status codes}
	\label{table.17}
	\tabularnewline
	\toprule\noalign{}
	Status code   & Description                                          \\
	\midrule\noalign{}
	\endfirsthead
	\toprule\noalign{}
	Status code   & Description                                          \\
	\midrule\noalign{}
	\endhead
	\bottomrule\noalign{}
	\endlastfoot
	0x20\emph{nn} & \emph{nn}: the number of \emph{remaining metrics} to 
	read \\
\end{longtable}

\paragraph{Metrics list operation}
\label{metrics_list_operation}

\textbf{Generic Command}, \textbf{Generic Status} and \textbf{Generic
Data} work together to load Metrics records from the device. The
scenario is as shown in figure \ref{illustration.28}.

\begin{figure}[H]
	\caption{Metrics scenario}
	\label{illustration.28}
	\centering
	\includegraphics[scale=0.5]{Metrics.pdf}
\end{figure}

\begin{enumerate}
	
	\item
	      Client enables \textbf{Generic Status} notifications.
	\item
	      Client writes the \textbf{Generic 	Command}{[}Metrics:UPDATE\_STATUS(0x2002){]} 
	      for the first record.
	\item
	      Client is notified of a \textbf{Generic Status} change, containing the new status, 
	      indicating 0x10\emph{nn-\/-}, with \emph{nn} the number of \textbf{Generic Data} 
	      records available. If \emph{nn} \textgreater{} 0, the scenario continues at point 4. 
	      Otherwise it is complete.
	\item
	      Client reads the \textbf{Generic Data} obtaining the data of the record
	\item
	      Client writes the \textbf{Generic Command}{[}Metrics:NEXT(0x2001){]}
	      for the next record. 
	      Note that probably this also indicates that the record is successfully stored by the Client, 
	      indicating that the Server can delete the record. 
	      So if you still need to sync with Nexxtmove.me, 
           you should not perform above scenario with another app.
	\item
	      Client is notified of a \textbf{Generic Status} change, containing the new status, 
	      indicating 0x10\emph{nn-\/-}, with \emph{nn} the number of
	      \textbf{Generic Data} records available. 
	      If \emph{nn} \textgreater{} 0, the scenario continues at point 4.
	      Otherwise it is complete.
	\item
	      Client disables \textbf{Generic Status} notifications.
\end{enumerate}

\footnote{Can this be confirmed}

The number of Metrics records that can be stored by the \emph{Nexxtender Home} is limited. 
The exact limit is not known. 
If the limit is reached, it is expected that older records are overwritten.

\paragraph{Generic Data format}
\label{generic_data_format_1}

Returns 20 bytes related to an event. 
The 20 bytes are structured as shown in tables \ref{table.77} and \ref{table.76o}.

Until now, no Metrics messages where seen in BLE traces.

\begin{longtable}{|l|l|l|l|l|l|l|l|l|l|l|l|l|l|l|l|l|l|l|l|}
	\caption{Generic Data data format overview}
	\label{table.77}\\  %<---- Problem
	\hline
	0 & 1 & 2 & 3 & 4 & 5 & 6 & 7 & 8 & 9 & 10 & 11 & 12 & 13 & 14 & 15 & 16 
	& 17 & 18 & 19 \\
	\hline
	\multicolumn{20}{|l|}{unknown} \\
	\hline
\end{longtable}

\hfill\break

\begin{longtable}[]{@{}lllll@{}}
	\caption{Generic Data data format detail}
	\label{table.76o}\\  %<---- Problem
	\tabularnewline
	\toprule\noalign{}
	Offset & Length & Name    & Type    & Description \\
	\midrule\noalign{}
	\endfirsthead
	\toprule\noalign{}
	Offset & Length & Name    & Type    & Description \\
	\midrule\noalign{}
	\endhead
	\bottomrule\noalign{}
	\endlastfoot
	0      & 20     & unknown & unknown &             \\
\end{longtable}

\subsubsection{Badge operations}
\label{badge_operations}

\paragraph{Badge operation ids}
\label{badge_operation_ids}

Table \ref{table.18} lists the \emph{Badge operation} id\textquotesingle s.

\begin{longtable}[]{@{}ll@{}}
	\caption{Badge operation ids}
	\label{table.18}
	\tabularnewline
	\toprule\noalign{}
	Operation id & Description                        \\
	\midrule\noalign{}
	\endfirsthead
	\toprule\noalign{}
	Operation id & Description                        \\
	\midrule\noalign{}
	\endhead
	\bottomrule\noalign{}
	\endlastfoot
	0x3001       & Add badge in default charging mode \\
	0x3002       & Add badge in MAX charging mode     \\
	0x3004       & Delete badge                       \\
	0x3005       & Badge list start                   \\
	0x3006       & Badge list next                    \\
\end{longtable}

\paragraph{Badge status}
\label{badge_status}

Possible Badge\emph{status} values are shown in table \ref{table.19}.

\begin{longtable}[]{@{}ll@{}}
	\caption{Badge status codes}
	\label{table.19}
	\tabularnewline
	\toprule\noalign{}
	Status code & Description  \\
	\midrule\noalign{}
	\endfirsthead
	\toprule\noalign{}
	Status code & Description  \\
	\midrule\noalign{}
	\endhead
	\bottomrule\noalign{}
	\endlastfoot
	0x3001      & WAIT\_ADD    \\
	0x3002      & WAIT\_ADD    \\
	0x3004      & WAIT\_DELETE \\
	0x3005      & NEXT         \\
	0x3007      & FINISH       \\
	0x3008      & ADDED        \\
	0x3009      & EXISTS       \\
	Other       & UNKNOWN      \\
\end{longtable}

\paragraph{Add badge in default charging mode operation}
\label{add_badge_in_default_charging_mode_operation}

This operation starts the procedure for adding a badge to the whitelist 
of the \emph{Nexxtender Home}. 
The whitelist is the list of badges allowed in Private mode. 
The new badge will always load in the default charging mode; 
i.e. ECO or MAX as set in the \emph{Mode} field from section \secref{config_1_0_format}.

The procedure is as shown in figure \ref{illustration.21}.

\begin{figure}[H]
	\caption{Add badge in default charging mode operation}
	\label{illustration.21}
	\centering
	\includegraphics[scale=0.5]{BadgeAddDefault.pdf}
\end{figure}

\begin{enumerate}
	
	\item
	      Client writes \textbf{Generic Command}{[}Badge:Add badge in default
	      charging mode (0x3001){]}.
	\item
	      Client is notified of a \textbf{Generic Status} change, 
	      containing the new status code Badge:WAIT\_ADD(0x3001), 
	      indicating that the \emph{Nexxtender Home} is waiting for the new badge
	      to be presented by the user. 
	      After the badge is presented, the Client is notified with a \textbf{Generic Status} change.
	\item
	      If the returned status is Badge:ADDED(0x3008), the new badge is successfully added,
	\item
	      and the client reads \textbf{Generic Data} and gets the UID data of the new badge.
	\item
	      If the returned status is Badge:EXISTS(0x3009), 
	      the presented badge already existed in the \emph{Nexxtender Home}. 
	      Nothing is changed.
\end{enumerate}

The ISO/IEC 14443-3 UID of the added badge is returned in step 4,
preceded with 1 byte indicating its length, as shown in tables \ref{table.20},
ref{table.21} and \ref{table.22}.

\begin{longtable}[]{@{}lllll@{}}
	\caption{4-byte UID return value}
	\label{table.20}
	\tabularnewline
	\toprule\noalign{}
	0 & 1    & 2    & 3    & 4    \\
	\midrule\noalign{}
	\endfirsthead
	\toprule\noalign{}
	0 & 1    & 2    & 3    & 4    \\
	\midrule\noalign{}
	\endhead
	\bottomrule\noalign{}
	\endlastfoot
	4 & UID0 & UID1 & UID2 & UID3 \\
\end{longtable}

\hfill\break

\begin{longtable}[]{@{}llllllll@{}}
	\caption{7-byte UID return value}
	\label{table.21}
	\tabularnewline
	\toprule\noalign{}
	0 & 1    & 2    & 3    & 4    & 5    & 6    & 7    \\
	\midrule\noalign{}
	\endfirsthead
	\toprule\noalign{}
	0 & 1    & 2    & 3    & 4    & 5    & 6    & 7    \\
	\midrule\noalign{}
	\endhead
	\bottomrule\noalign{}
	\endlastfoot
	7 & UID0 & UID1 & UID2 & UID3 & UID4 & UID5 & UID6 \\
\end{longtable}

\hfill\break

\begin{longtable}[]{@{}lllllllllll@{}}
	\caption{10-byte UID return value}
	\label{table.22}
	\tabularnewline
	\toprule\noalign{}
	0  & 1    & 2    & 3    & 4    & 5    & 6    & 7    & 8    & 9    & 10 \\
	\midrule\noalign{}
	\endfirsthead
	\toprule\noalign{}
	0  & 1    & 2    & 3    & 4    & 5    & 6    & 7    & 8    & 9    & 10 \\
	\midrule\noalign{}
	\endhead
	\bottomrule\noalign{}
	\endlastfoot
	10 & UID0 & UID1 & UID2 & UID3 & UID4 & UID5 & UID6 & UID7 & UID8 &    
	UID9 \\
\end{longtable}

The 7-byte version is in use on my \emph{Nexxtender Home}. 
I don't know if the 4-byte or 10-byte format is accepted.

Note that at least for older \emph{Nexxtender Home} firmware versions (like 2.53.2), 
this scenario does not seem to work. 
Regardless of the Default charging settings, 
the newly added badge will always load in MAX mode.

\paragraph{Add badge in MAX charging mode operation}
\label{add_badge_in_max_charging_mode_operation}

This operation starts the procedure for adding a badge to the \emph{Nexxtender Home}. 
The new badge will always load in the MAX charging mode, 
independently of the default charging mode.

The procedure is exactly the same as for
\secref{add_badge_in_default_charging_mode_operation}, 
except that command {[}Badge:Add badge in default charging mode (0x3001){]} is replaced with
{[}Badge:Add badge in MAX charging mode (0x3002){]}. 
Also the status code Badge:WAIT\_ADD(0x3002) is expected instead of Badge:WAIT\_ADD(0x3001).
See figure \ref{illustration.22}.

\begin{figure}[H]
	\caption{Add badge in MAX charging mode operation}
	\label{illustration.22}
	\centering
	\includegraphics[scale=0.5]{BadgeAddMAX.pdf}
\end{figure}

\hfill\break

\paragraph{Delete Badge operation}
\label{delete_badge_operation}

This operation starts the procedure for deleting a badge from the badge
whitelist of the \emph{Nexxtender Home}.

The procedure is as shown in figure \ref{illustration.23}.

\begin{figure}[H]
	\caption{Delete Badge operation}
	\label{illustration.23}
	\centering
	\includegraphics[scale=0.5]{BadgeDelete.pdf}
\end{figure}

\begin{enumerate}
	\item
	      Client writes \textbf{Generic Command}{[}Badge:Delete badge (0x3004){]}
	\item
	      Client is notified of a \textbf{Generic Status} change, containing the
	      Badge:WAIT\_DELETE(0x3004) status.
	\item
	      The client writes a \textbf{Generic Data} with the UID of the badge that must be deleted. 
	      The format of the UID equals the one described in \secref{add_badge_in_default_charging_mode_operation}.
	\item
	      Client is notified of a \textbf{Generic Status} change, containing the
	      Badge:FINISH(0x3004) status, 
	      indicating that the badge was successfully deleted\footnote{Are we sure of this? Do we need to wait for a notification on
	      	\textbf{Generic Status} with a status code?}.
\end{enumerate}

\paragraph{Badge list operation}
\label{badge_list_operation}

This operation starts the procedure for obtaining the badge whitelist of the \emph{Nexxtender Home}.

The procedure is as shown in figure \ref{illustration.24}.

\begin{figure}[H]
	\caption{Badge list operation}
	\label{illustration.24}
	\centering
	\includegraphics[scale=0.5]{BadgeList.pdf}
\end{figure}

\begin{enumerate}
	
	\item
	      Client writes \textbf{Generic Command}{[}Badge:Badge list start (0x3005){]}
	\item
	      Client is notified of a \textbf{Generic Status} change, containing the new status.
	\item
	      If the status is Badge:NEXT(0x3005), 
 	      then the client reads \textbf{Generic Data} and gets the UID of the badge. 
	      The returned format equals the one described in
	      \secref{add_badge_in_default_charging_mode_operation}. 
	      The context (Default or MAX charging mode) is apparently not returned,
	\item
	      and the Client writes \textbf{Generic Command}{[}Badge:Badge list next (0x3006){]},
	\item
	      and the Client is notified of a \textbf{Generic Status} change,
	      containing the new status. 
	     Repeat from step 3.
	\item
	      If the status is Badge:FINISH(0x3007), the Client has successfully
	      obtained the complete list.
	\item
	      If the status is anything else, there was an error.
\end{enumerate}

\subsubsection{Time operation}
\label{time_operation}

\paragraph{Time operations ids}
\label{time_operations_ids}

Table \ref{table.23} lists the \emph{Time operation} id\textquotesingle s.

\begin{longtable}[]{@{}ll@{}}
	\caption{Time operation ids}
	\label{table.23}
	\tabularnewline
	\toprule\noalign{}
	Operation Id & Description \\
	\midrule\noalign{}
	\endfirsthead
	\toprule\noalign{}
	Operation Id & Description \\
	\midrule\noalign{}
	\endhead
	\bottomrule\noalign{}
	\endlastfoot
	0x4001       & Time Set    \\
	0x4002       & Time Get    \\
\end{longtable}

\paragraph{Time status}
\label{time_status}

Possible Time \emph{status} values are shown in table \ref{table.24}.

\begin{longtable}[]{@{}ll@{}}
	\caption{Time status codes}
	\label{table.24}
	\tabularnewline
	\toprule\noalign{}
	Status code & Description \\
	\midrule\noalign{}
	\endfirsthead
	\toprule\noalign{}
	Status code & Description \\
	\midrule\noalign{}
	\endhead
	\bottomrule\noalign{}
	\endlastfoot
	0x4001      & READY       \\
	0x4002      & SUCCESS     \\
	0x4003      & POPPED      \\
	Other       & UNKNOWN     \\
\end{longtable}

\paragraph{Time Set operation}
\label{time_set_operation}

This is an operation that writes the time.
See figure \ref{illustration.29}.

\begin{figure}[H]
	\caption{Time Set Operation}
	\label{illustration.29}
	\centering
	\includegraphics[scale=0.5]{TimeSet.pdf}
\end{figure}

\begin{enumerate}
	
	\item
	      Client writes \textbf{Generic Command}{[}Time:Time Set(0x4001){]}
	\item
	      Client is notified of a \textbf{Generic Status} change, 
	      containing the new status code Time:READY(0x4001), 
	      indicating that the generic data is ready to be written.
	\item
	      The client writes a \textbf{Generic Data} with the time to set.
	\item
	      Client is notified of a \textbf{Generic Status} change, 
	      containing the new Time:SUCCESS(0x4002), 
	      indicating that the time was successfully written. 
	      But in practice it seems that the \emph{Nexxtender Home}
	      never sends this Time:SUCCESS(0x4002).
\end{enumerate}

The data format for the \textbf{Generic Data} is shown in tables \ref{table.25} and \ref{table.26}.

\begin{longtable}[]{|l|l|l|l|}
	\caption{Time Set data format overview}
	\label{table.25}\\  %<---- Problem
	\hline
	0 & 1 & 2 & 3 \\
	\hline
	\multicolumn{4}{|l|}{timestamp} \\
	\hline
\end{longtable}

\hfill\break

\begin{longtable}[]{@{}lll@{}}
	\caption{Time Set data format detail}
	\label{table.26}
	\tabularnewline
	\toprule\noalign{}
	Field     & Bytes & Description         \\
	\midrule\noalign{}
	\endfirsthead
	\toprule\noalign{}
	Field     & Bytes & Description         \\
	\midrule\noalign{}
	\endhead
	\bottomrule\noalign{}
	\endlastfoot
	timestamp & 0:3   & Current time in UTC 
	\href{https://en.wikipedia.org/wiki/Unix_time}{Unix Time}. 
	4 byte unsigned integer in
	\href{https://en.wikipedia.org/wiki/Endianness}{Little-endian} \\
\end{longtable}

\paragraph{Time Get operation}
\label{time_get_operation}

This is an operation that reads the time.
See figure \ref{illustration.10}.

\begin{figure}[H]
	\caption{Time Get Operation}
	\label{illustration.10}
	\centering
	\includegraphics[scale=0.5]{TimeGet.pdf}
\end{figure}

\begin{enumerate}
	
	\item
	      Client writes \textbf{Generic Command}{[}Time:Time Get(0x4002){]}
	\item
	      Client is notified of a \textbf{Generic Status} change, 
	      containing the new status Time:POPPED(0x4003), 
	      indicating that the generic data is ready to be read.
	\item
	      The client reads \textbf{Generic Data} and gets the requested time.
\end{enumerate}

The data format for the \textbf{Generic Data} is as in section \secref{time_set_operation}.

\subsubsection{Config operations}
\label{config_operations}

These operations allow to get and set the \emph{Nexxtender Home}
\textquotesingle s configurable parameters. 
3 different formats exist to represent the configuration data. 
Which one to use depends on the \emph{Nexxtender Home} charger\textquotesingle s 
Firmware Revision String (according to table \secref{table.8}).

Table \ref{table.27} shows the relation between firmware revision and configuration data format.

\begin{longtable}[]{@{}ll@{}}
	\caption{Configuration data format to use}
	\label{table.27}
	\tabularnewline
	\toprule\noalign{}
	Firmware Revision                                         & Configuration data format \\
	\midrule\noalign{}
	\endfirsthead
	\toprule\noalign{}
	Firmware Revision                                         & Configuration data format \\
	\midrule\noalign{}
	\endhead
	\bottomrule\noalign{}
	\endlastfoot
	Firmware Revision \textless{} 1.1.0                       & Config\_1\_0              \\
	1.1.0 \ensuremath{\le} Firmware Revision \textless{} 3.50 & Config\_1\_1              \\
	Firmware Revision \ensuremath{\ge} 3.50                   & Config\_CBOR              \\
\end{longtable}

\paragraph{Config\_1\_0 format}
\label{config_1_0_format}

In Config\_1\_0, 13 configuration bytes are exchanged.
The 13 bytes are structured as shown in tables \ref{table.28} and \ref{table.31}.

\begin{longtable}[]{|l|l|l|l|l|l|l|l|l|l|l|l|l|}
	\caption{Config\_1\_0 format overview}
	\label{table.28}\\  %<---- Problem
	\hline
	0 & 1 & 2 & 3 & 4 & 5 & 6 & 7 & 8 & 9 & 10 & 11 & 12 \\
	\hline
	maxGrid & mode & safe & \multicolumn{2}{l|}{touWeekStart} & 
	\multicolumn{2}{l|}{touWeekEnd} & 
	\multicolumn{2}{l|}{touWeekendStart} & 
	\multicolumn{2}{l|}{touWeekendEnd} & 
	\multicolumn{2}{l|}{crc16} \\
	\hline
\end{longtable}

\begin{longtable}[]{@{}llllp{9.5cm}@{}}
	\caption{Config \_1\_0 data format detail}
	\label{table.31}
	\tabularnewline
	\toprule\noalign{}
	Offset & Length & Name            & Type   & Description                                    \\
	\midrule\noalign{}
	\endfirsthead
	\toprule\noalign{}
	Offset & Length & Name            & Type   & Description                                    \\
	\midrule\noalign{}
	\endhead
	\bottomrule\noalign{}
	\endlastfoot
	0      & 1      & maxGrid         & uint8  & Maximum allowed grid consumption limit in A.   
	Set this equal or smaller to the total A values of the main grid fuses.
	If it is 3 * 16A, set it to 48A. 
	If it is 1 * 50A, set it to 50A. 
	The \emph{Nexxtender Home} will make sure that total grid consumption will
	remain less than the specified value. 
	You can set this to a lower value	if you have a peak tariff contract. \\
	1      & 1      & mode            & uint8  & The default charging mode. Codes ECO\_PRIVATE, 
	MAX\_PRIVATE, ECO\_OPEN, MAX\_OPEN according to table
	\secref{table.30}. \\
	2      & 1      & safe            & uint8  & Minimum charging speed in A for the device.    
	Certified chargers are required to provide a minimum of 6A, 
	so the normal \emph{safe} value is 6A. \\
	3      & 2      & touWeekStart    & int16  & Off peak charging start time of each weekday. 
	Coded in minutes since midnight. \\
	5      & 2      & touWeekEnd      & int16  & Off peak charging end time of each weekday.    
	Coded in minutes since midnight. \\
	7      & 2      & touWeekendStart & int16  & Off peak charging start time of each weekend day. 
	Coded in minutes since midnight. \\
	9      & 2      & touWeekendEnd   & int16  & Off peak charging end time of each	weekend day. 
	Coded in minutes since midnight. \\
	11     & 2      & crc16           & uint16 & CRC-16/MODBUS over previous 11 bytes           \\
\end{longtable}

The following table shows the coding for the \emph{mode} field. 
b0 codes if charging is in ECO or MAX, b2 codes if charging is PRIVATE or OPEN.

\begin{longtable}[]{@{}llllllllllp{6cm}@{}}
	\caption{mode coding}
	\label{table.30}\\  %<---- Problem
	\tabularnewline
	\toprule\noalign{}
	b7 & b6 & b5 & b4 & b3 & b2 & b1 & b0 & Value & Name         & Description            \\
	\midrule\noalign{}
	\endfirsthead
	\toprule\noalign{}
	b7 & b6 & b5 & b4 & b3 & b2 & b1 & b0 & Value & Name         & Description            \\
	\midrule\noalign{}
	\endhead
	\bottomrule\noalign{}
	\endlastfoot
	0  & 0  & 0  & 0  & 0  & 0  & 0  & 0  & 0     & ECO\_PRIVATE & ECO mode charging      
	with a whitelisted badge \\
	0  & 0  & 0  & 0  & 0  & 0  & 0  & 1  & 1     & MAX\_PRIVATE & MAX mode charging      
	with a whitelisted badge \\
	0  & 0  & 0  & 0  & 0  & 1  & 0  & 0  & 4     & ECO\_OPEN    & ECO mode charging with 
	any badge \\
	0  & 0  & 0  & 0  & 0  & 1  & 0  & 1  & 5     & MAX\_OPEN    & MAX mode charging with 
	any badge \\
	x  & x  & x  & x  & x  & x  & x  & x  & x     & RFU          & Any other value: RFU   \\
\end{longtable}

\paragraph{Config\_1\_1 format}
\label{config_1_1_format}

In Config\_1\_1, 15 configuration bytes are exchanged. 
The 15 bytes are structured as shown in tables \ref{table.32} and \ref{table.78}.
The fields maxDevice and networkType are new compared to the ones from 
section \secref{config_1_0_format}.

\begin{longtable}{|l|l|l|l|l|l|l|l|l|l|l|l|l|l|l|}
	\caption{Config\_1\_1 format overview}
	\label{table.32}\\  %<---- Problem
	\hline
	0 & 1 & 2 & 3 & 4 & 5 & 6 & 7 & 8 & 9 & 10 & 11 & 12 & 13 & 14 \\
	\hline
	\scriptsize maxGrid & \scriptsize maxDevice & \scriptsize mode & \scriptsize safe & \scriptsize networkType & \multicolumn{2}{l}{%
	\scriptsize touWeekStart} & \multicolumn{2}{|l}{%
	\scriptsize touWeekEnd} & \multicolumn{2}{|l}{%
	\scriptsize touWeekendStart} & \multicolumn{2}{|l}{%
	\scriptsize touWeekendEnd} & \multicolumn{2}{|l|}{%
	\scriptsize crc16} \\
	\hline
\end{longtable}

\begin{longtable}[]{@{}llllp{9cm}@{}}
	\caption{Config \_1\_1 data format detail}
	\label{table.78}
	\tabularnewline
	\toprule\noalign{}
	Offset & Length & Name            & Type   & Description                                 \\
	\midrule\noalign{}
	\endfirsthead
	\toprule\noalign{}
	Offset & Length & Name            & Type   & Description                                 \\
	\midrule\noalign{}
	\endhead
	\bottomrule\noalign{}
	\endlastfoot
	0      & 1      & maxGrid         & uint8  & See \secref{config_1_0_format}              \\
	1      & 1      & maxDevice       & uint8  & Maximum allowed charging speed in A for the 
	EV device. \emph{safe} \textless= \emph{maxDevice} \textless=
	\emph{maxGrid}. \\
	2      & 1      & mode            & uint8  & See \secref{config_1_0_format}              \\
	3      & 1      & safe            & uint8  & See \secref{config_1_0_format}              \\
	4      & 1      & networkType     & uint8  & Codes the network type according to tables  
	\secref{table.33}  and \secref{table.75}\\
	5      & 2      & touWeekStart    & int16  & See                                         \\
	7      & 2      & touWeekEnd      & int16  & See                                         
	\secref{config_1_0_format} \\
	9      & 2      & touWeekendStart & int16  & See                                         
	\secref{config_1_0_format} \\
	11     & 2      & touWeekendEnd   & int16  & See                                         
	\secref{config_1_0_format} \\
	13     & 2      & crc16           & uint16 & CRC-16/MODBUS over previous 13 bytes        \\
\end{longtable}

Tables \ref{table.33} and \ref{table.75} show the coding for the \emph{networkType} field.

\begin{longtable}[]{@{}lllllllllll@{}}
	\caption{networkType coding in Config Get}
	\label{table.33}
	\tabularnewline
	\toprule\noalign{}
	b7 & b6 & b5 & b4 & b3 & b2 & b1 & b0 & Value  & Name       & Description                 \\
	\midrule\noalign{}
	\endfirsthead
	\toprule\noalign{}
	b7 & b6 & b5 & b4 & b3 & b2 & b1 & b0 & Value  & Name       & Description                 \\
	\midrule\noalign{}
	\endhead
	\bottomrule\noalign{}
	\endlastfoot
	0  & 0  & 0  & 0  & 0  & 0  & 0  & x  & 0 or 1 & Mono/Tri+N & Single-phase or             
	Three-phase with neutral \\
	0  & 0  & 0  & 0  & 0  & 0  & 1  & x  & 2 or 3 & Tri        & Three-phase without neutral \\
	x  & x  & x  & x  & x  & x  & x  & x  & x      & RFU        & Any other value: RFU        \\
\end{longtable}

\begin{longtable}[]{@{}lllllllllll@{}}
	\caption{networkType coding in Config Set}
	\label{table.75}
	\tabularnewline
	\toprule\noalign{}
	b7 & b6 & b5 & b4 & b3 & b2 & b1 & b0 & Value & Name       & Description                 \\
	\midrule\noalign{}
	\endfirsthead
	\toprule\noalign{}
	b7 & b6 & b5 & b4 & b3 & b2 & b1 & b0 & Value & Name       & Description                 \\
	\midrule\noalign{}
	\endhead
	\bottomrule\noalign{}
	\endlastfoot
	0  & 0  & 0  & 0  & 0  & 0  & 0  & 0  & 0     & Mono/Tri+N & Single-phase or             
	Three-phase with neutral \\
	0  & 0  & 0  & 0  & 0  & 0  & 1  & 1  & 3     & Tri        & Three-phase without neutral \\
	x  & x  & x  & x  & x  & x  & x  & x  & x     & RFU        & Any other value: RFU        \\
\end{longtable}
\paragraph{Config\_CBOR format}
\label{config_cbor_format}

In Config\_CBOR, a CBOR coding is used. 
The top level is a CBOR Map containing 2 \emph{Key}/\emph{Value} pairs:

\begin{itemize}
	
	\item
	      Pair 1 (Not sure what this pair is used for)
	      
	      \begin{itemize}
	      	  
	      	\item
	      	      \emph{Key}: CBOR Integer 0
	      	\item
	      	      \emph{Value}: CBOR Map containing 2 \emph{Key}/\emph{Value} pairs
	      	      
	      	      \begin{itemize}
	      	      	    
	      	      	\item
	      	      	      Pair 1\footnote{What is this?}
	      	      	      
	      	      	      \begin{itemize}
	      	      	      	      
	      	      	      	\item
	      	      	      	      \emph{Key}: CBOR Integer 1
	      	      	      	\item
	      	      	      	      \emph{Value}: CBOR Integer 1
	      	      	      \end{itemize}
	      	      	\item
	      	      	      Pair 2\footnote{What is this?}
	      	      	      
	      	      	      \begin{itemize}
	      	      	      	      
	      	      	      	\item
	      	      	      	      \emph{Key}: CBOR Integer 2
	      	      	      	\item
	      	      	      	      \emph{Value}: CBOR Integer 1
	      	      	      \end{itemize}
	      	      \end{itemize}
	      \end{itemize}
	\item
	      Pair 2
	      
	      \begin{itemize}
	      	  
	      	\item
	      	      \emph{Key}: CBOR Integer 1
	      	\item
	      	      \emph{Value}: CBOR Map containing the Nexxtender configuration values. 
	               Each pair represents a configuration value. 
	               All \emph{Keys} and \emph{Values} in this map are CBOR Integers. 
	               The possible \emph{Key} values are shown in table \ref{table.34}.
	      \end{itemize}
\end{itemize}

In addition, 
a CRC-16/MODBUS crc is computed over the complete CBOR coded message and appended to it.
Table \ref{table.34} shows the structure of the data.

The keys \emph{iEvseMin} and \emph{iCapacity} are new compared to the
ones from \secref{config_1_1_format}.

\begin{longtable}[]{@{}lllp{12cm}@{}}
	\caption{Config CBOR Map format}
	\label{table.34}
	\tabularnewline
	\toprule\noalign{}
	\multicolumn{2}{@{}l}{%
	Key} & Value/Value name & Description \\
	\midrule\noalign{}
	\endfirsthead
	\toprule\noalign{}
	\multicolumn{2}{@{}l}{%
	Key} & Value/Value name & Description \\
	\midrule\noalign{}
	\endhead
	\bottomrule\noalign{}
	\endlastfoot
	0 &    &                    & Map                                                                                 \\
	  & 1  & 1                  & ?                                                                                   \\
	  & 2  & 1                  & ?                                                                                   \\
	1 &    &                    & Map\footnote{The possible key values shown in the next rows for this map are        
	ordered in increasing integer value. 
	That is not a requirement.
	Actually the Nexxtmove app creates CBOR messages in which the pairs	are not sorted according to key value.}\\
	  & 1  & chargeMode         & See \emph{mode} in                                                                  
	\secref{config_1_1_format} \\
	  & 2  & modbusSlaveAddress & Not used                                                                            \\
	  & 3  & cycleRate          & Not used                                                                            \\
	  & 4  & iMax               & See \emph{maxGrid} in                                                               
	\secref{config_1_1_format} \\
	  & 5  & iEvseMax           & See \emph{maxDevice} in                                                             
	\secref{config_1_1_format} \\
	  & 6  & iEvseMin           & Not used\footnote{Nexxtender home stores it in Config\_CBOR.minDevice, but does not 
	use it. 
	It is also unclear what the difference with iLevel1 would be.} \\
	  & 7  & iLevel1            & See \emph{safe}in                                                                   
	\secref{config_1_1_format} \\
	  & 8  & solarMode          & Not used                                                                            \\
	  & 9  & phaseSeq           & Only b1 is used as b0 in \emph{networkType}, see                                    
	\secref{config_1_1_format} \\
	  & 10 & chargingPhases     & Not used                                                                            \\
	  & 11 & blePin             & Not used                                                                            \\
	  & 12 & touWeekStart       & See \emph{touWeekStart} in                                                          
	\secref{config_1_1_format} \\
	  & 13 & touWeekStop        & See \emph{touWeekEnd} in                                                            
	\secref{config_1_1_format} \\
	  & 14 & touWeekendStart    & See \emph{touWeekendStart} in                                                       
	\secref{config_1_1_format} \\
	  & 15 & touWeekendStop     & See \emph{touWeekendEnd} in                                                         
	\secref{config_1_1_format}\\
	  & 16 & timezone           & Not used                                                                            \\
	  & 17 & relayOffPeriod     & Not used                                                                            \\
	  & 18 & externalRegulation & Not used                                                                            \\
	  & 19 & iCapacity          & Peak grid current limit\footnote{Is this for 1 phase or all phases together?}.      
	\emph{iLevel1} \textless= \emph{iCapacity} \textless= \emph{iMax}. You
	can set this to a lower value if you have a peak tariff contract. In
	Config\_CBOR format it is not needed to use \emph{iMax} (\emph{maxGrid})
	for that purpose.
	\secref{config_1_0_format} and
	\secref{config_1_1_format} behave as if \emph{iCapacity} is set to \emph{maxGrid}. \\
\end{longtable}

\paragraph{Config operations ids}
\label{config_operations_ids}

Table \ref{table.35} lists the \emph{Config Operation} id\textquotesingle s.
The {[}Config:Config
	Set(0x5001){]} and {[}Config:Config
Get(0x5002){]} operations are used to exchange the
\secref{config_1_0_format} and
\secref{config_1_0_format}. The {[}Config:Config CBOR Set(0x5003){]} and {[}Config:Config
CBOR Get(0x5004){]} operations are used to exchange the
\secref{config_cbor_format}.

\begin{longtable}[]{@{}ll@{}}
	\caption{Config operation ids}
	\label{table.35}
	\tabularnewline
	\toprule\noalign{}
	Operation id & Description     \\
	\midrule\noalign{}
	\endfirsthead
	\toprule\noalign{}
	Operation id & Description     \\
	\midrule\noalign{}
	\endhead
	\bottomrule\noalign{}
	\endlastfoot
	0x5001       & Config Set      \\
	0x5002       & Config Get      \\
	0x5003       & Config CBOR Set \\
	0x5004       & Config CBOR Get \\
\end{longtable}

\paragraph{Config status}
\label{config_status}

Possible Config \emph{status} values are shown in table \ref{table.36}.

\begin{longtable}[]{@{}ll@{}}
	\caption{Config status codes}
	\label{table.36}
	\tabularnewline
	\toprule\noalign{}
	Config status & Description                       \\
	\midrule\noalign{}
	\endfirsthead
	\toprule\noalign{}
	Config status & Description                       \\
	\midrule\noalign{}
	\endhead
	\bottomrule\noalign{}
	\endlastfoot
	0x5001        & READY (after a Config Set)        \\
	0x5002        & SUCCESS (after a Config Set)      \\
	0x5003        & POPPED (after a Config Get)       \\
	0x5004        & READY (after a Config CBOR Set)   \\
	0x5005        & SUCCESS (after a Config CBOR Set) \\
	0x5006        & POPPED (after a Config CBOR Get)  \\
	Other         & UNKNOWN                           \\
\end{longtable}

\paragraph{Config Set operation}
\label{config_set_operation}

This is an operation that writes several configuration values. 
The expected success scenario is as shown in figure \ref{illustration.4}.

\begin{figure}[H]
	\caption{Config Set Operation}
	\label{illustration.4}
	\centering
	\includegraphics[scale=0.5]{SetConfigSuccess.pdf}
\end{figure}

\begin{enumerate}
	
	\item
	      Client writes \textbf{Generic Command}{[}Config:Config Set(0x5001){]}
	\item
	      Client is notified of a \textbf{Generic Status} change, 
	      containing the status code Config:READY(0x5001), 
	      indicating that the generic data is ready to be written.
	\item
	      The client writes a \textbf{Generic Data} with the actual data that
	      must be written for Config:Config Set(0x5001). 
	      This is either the \secref{config_1_0_format} or \secref{config_1_0_format}.
	\item
	      Client is notified of a \textbf{Generic Status} change, 
	      containing the new status Config:SUCCESS(0x5002), 
	      indicating that the data was successfully written.
\end{enumerate}

\paragraph{Config CBOR Set operation}
\label{config_cbor_set_operation}

This is an operation that writes several configuration values. 
It is similar to \secref{config_set_operation}, but the input is in CBOR format. 
It is only available on \emph{Nexxtender Home} chargers with a Firmware Revision String of at least 3.50. 
The expected success scenario is as shown in figure \ref{illustration.5}.

\begin{figure}[H]
	\caption{Config CBOR Set Operation}
	\label{illustration.5}
	\centering
	\includegraphics[scale=0.5]{SetConfigCBORSuccess.pdf}
\end{figure}

\begin{enumerate}
	
	\item
	      Client writes \textbf{Generic Command}{[}Config:Config CBOR
	      Set(0x5003){]}
	\item
	      Client is notified of a \textbf{Generic Status} change, 
	      containing the new status Config:READY(0x5004), 
	      indicating that the generic data is ready to be written.
	\item
	      The client writes a \textbf{Generic Data} with the actual data that
	      must be written for Config Set. 
	      This is the \secref{config_cbor_format}.
	\item
	      Client is notified of a \textbf{Generic Status} change, 
	      containing the new status Config:SUCCESS(0x5005), 
	      indicating that the data was successfully written.
\end{enumerate}

\paragraph{Config Get Operation}
\label{config_get_operation}

This is an operation that reads several configuration values. 
The expected success scenario is as shown in figure \ref{illustration.3}.

\begin{figure}[H]
	\caption{Config Get Operation}
	\label{illustration.3}
	\centering
	\includegraphics[scale=0.5]{GetConfigSuccess.pdf}
\end{figure}

\begin{enumerate}
	
	\item
	      Client writes \textbf{Generic Command}{[}Config:Config Get(0x5002){]}
	\item
	      Client is notified of a \textbf{Generic Status} change, 
	      containing the new status Config:POPPED(0x5003), 
	      indicating that the generic data is ready to be read.
	\item
	      The client reads \textbf{Generic Data} and gets the requested data.
	      This is either the \secref{config_1_0_format} or \secref{config_1_1_format}.
\end{enumerate}

\paragraph{Config CBOR Get	operation}
\label{config_cbor_get_operation}

This is an operation that reads several configuration values. 
The expected success scenario is as shown in figure \ref{illustration.6}.

\begin{figure}[H]
	\caption{Config CBOR Get Operation}
	\label{illustration.6}
	\centering
	\includegraphics[scale=0.5]{GetConfigCBORSuccess.pdf}
\end{figure}

\begin{enumerate}
	
	\item
	      The client writes \textbf{Generic Command}{[}Config:Config CBOR Get(0x5004){]}
	\item
	      Client is notified of a \textbf{Generic Status} change, 
	      containing the new status Config:POPPED(0x5006), 
	      indicating that the generic data is ready to be read.
	\item
	      The client reads \textbf{Generic Data} and gets the requested data.
	      This is the \secref{config_cbor_format}.
\end{enumerate}

\subsection{CDR characteristics}
\label{cdr_characteristics}

Table \ref{table.37} shows which characteristics of the \emph{CDR} service (also known as
\emph{Transaction Service}) are used.

\textbf{Note}: If the charger is an \emph{Nexxtender Home}, 
these characteristics are send under the Generic (C1) service. 
If the charger is an \emph{Nexxtender Mobile}, these characteristics are send under the CDR(C1) service.

\begin{longtable}[]{@{}llllllll@{}}
	\caption{GATT CDR service characteristics used}
	\label{table.37}
	\tabularnewline
	\toprule\noalign{}
	Handle & UUID & Name & R & W & WWR & N & Coding \\
	\midrule\noalign{}
	\endfirsthead
	\toprule\noalign{}
	Handle & UUID & Name & R & W & WWR & N & Coding \\
	\midrule\noalign{}
	\endhead
	\bottomrule\noalign{}
	\endlastfoot
	0x001b & fd47416a-95fb-4206-88b5-b4a8045f75\textbf{c2} & \textbf{CDR
	Command} & & P? & P? & & HEX \\
	0x001c & fd47416a-95fb-4206-88b5-b4a8045f75\textbf{c3} & \textbf{CDR
	Status} & P & & & N & HEX \\
	0x0020 & fd47416a-95fb-4206-88b5-b4a8045f75\textbf{c4} & \textbf{CDR
	Record} & P & & & & HEX \\
\end{longtable}

\textbf{CDR Command}, \textbf{CDR Status} and \textbf{CDR Record} work
together to load transaction records from the device. 
The scenario is as shown in figure \ref{illustration.18}.

\begin{figure}[H]
	\caption{CDR scenario}
	\label{illustration.18}
	\centering
	\includegraphics[scale=0.5]{CDR.pdf}
\end{figure}

\begin{enumerate}
	
	\item
	      Client enables \textbf{CDR Status} notifications.
	\item
	      Client reads the \textbf{CDR Status} obtaining the new status,
	      indicating \emph{nnnnnnnn}, the number of CDR records available. 
	      If \emph{nnnnnnnn} == 0, the scenario continues at point 6.
	\item
	      Client reads the \textbf{CDR Record} obtaining the data of the record.
	\item
	      Client writes the \textbf{CDR Command}{[}0x01{]} for the next record.
	      Note that probably this also indicates that the record is successfully stored by the Client, 
	      indicating that the Server can delete the record. 
	      So if you still need to sync with Nexxtmove.me, 
	      you should not perform above scenario with another app.
	\item
	      Client is notified of a \textbf{CDR Status} change, containing the new status, 
	      indicating \emph{nnnnnnnn-\/-}, the new number of CDR records available. 
          If \emph{nnnnnnnn} \textgreater{} 0, the scenario continues at point 3. 
	     Otherwise it is complete.
	\item
	      Client disables \textbf{CDR Status} notifications.
\end{enumerate}

The number of CDR records that can be stored by the Nexxtender Chargeris limited. 
The exact limit is not known. 
If the limit is reached, it isexpected that older records are overwritten.

\subsubsection{CDR Command operation ids}
\label{cdr_command_operation_ids}

Table \ref{table.55} lists the \emph{CDR Command} operation id\textquotesingle s.

\begin{longtable}[]{@{}lp{14cm}@{}}
	\caption{CDR Command operation ids}
	\label{table.55}
	\tabularnewline
	\toprule\noalign{}
	Operation Id & Description                                                    \\
	\midrule\noalign{}
	\endfirsthead
	\toprule\noalign{}
	Operation Id & Description                                                    \\
	\midrule\noalign{}
	\endhead
	\bottomrule\noalign{}
	\endlastfoot
	0x01         & Read next record? Probably also means that the previous record 
	read from \textbf{CDR Record} can be deleted from the Server. \\
	others?      &                                                                \\
\end{longtable}

\subsubsection{CDR Status status}
\label{cdr_status_status}

Possible \emph{CDR Status} values are shown in table \ref{table.56}.

\begin{longtable}[]{@{}ll@{}}
	\caption{CDR Status status codes}
	\label{table.56}
	\tabularnewline
	\toprule\noalign{}
	Status Code       & Description                                     \\
	\midrule\noalign{}
	\endfirsthead
	\toprule\noalign{}
	Status Code       & Description                                     \\
	\midrule\noalign{}
	\endhead
	\bottomrule\noalign{}
	\endlastfoot
	0x\emph{nnnnnnnn} & \emph{nnnnnnnn} is an uint32 with the number of 
	records still to read \\
\end{longtable}

\subsubsection{CDR Record format}
\label{cdr_record_format}

Returns a 32 bytes CDR record that summarizes a charging session. 
The 32 bytes are structured as shwon in tables \ref{table.53a} , \ref{table.53b} and \ref{table.54}.

\begin{longtable}[]{@{}|l|l|l|l|l|l|l|l|l|l|l|l|l|l|l|l|@{}}
	\caption{CDR Record data format overview, bytes 0 to 15}
	\label{table.53a}\\  %<---- Problem
	\hline
	0 & 1 & 2 & 3 & 4 & 5 & 6 & 7 & 8 & 9 & 10 & 11 & 12 & 13 & 14 & 15 \\
	\hline
	\multicolumn{4}{@{}|l}{%
	unknown1} & \multicolumn{4}{|l}{%
	sessionStartTime} & \multicolumn{4}{|l}{%
	sessionStartEnergy} & \multicolumn{4}{|l|}{%
	unknown2} \\
	\hline
\end{longtable}

\hfill\break

\begin{longtable}[]{@{}|l|l|l|l|l|l|l|l|l|l|l|l|l|l|l|l|@{}}
	\caption{CDR Record data format overview, bytes 16 to 31}
	\label{table.53b}\\  %<---- Problem
	\hline
	16
	  & 17 & 18 & 19 & 20 & 21 & 22 & 23 & 24 & 25 & 26 & 27 & 28 & 29 & 30 &   
	31 \\
	\hline
	\multicolumn{4}{|l}{%
	unknown3} & \multicolumn{4}{|l}{%
	sessionStopTime} & \multicolumn{4}{|l}{%
	sessionStopEnergy} & \multicolumn{2}{|l}{%
	unknown4} & \multicolumn{2}{|l@{}|}{%
	crc16} \\
	\hline
\end{longtable}

\hfill\break

\begin{longtable}[]{@{}llllp{9cm}@{}}
	\caption{CDR Record data format detail}
	\label{table.54}
	\tabularnewline
	\toprule\noalign{}
	Offset & Length & Name               & Type   & Description                          \\
	\midrule\noalign{}
	\endfirsthead
	\toprule\noalign{}
	Offset & Length & Name               & Type   & Description                          \\
	\midrule\noalign{}
	\endhead
	\bottomrule\noalign{}
	\endlastfoot
	0      & 4      & unknown1           & uint32 & Always 0?                            \\
	4      & 4      & sessionStartTime   & uint32 & Charging session start time          \\
	8      & 4      & sessionStartEnergy & uint32 & Total amount of energy in Wh since   
	the installation of the charger, at the moment charging starts. \\
	12     & 4      & unknown2           & uint32 & Always 0?                            \\
	16     & 4      & unknown3           & uint32 & Always 0?                            \\
	20     & 4      & sessionStopTime    & uint32 & Charging session stop time           \\
	24     & 4      & sessionStopEnergy  & uint32 & Total amount of energy in Wh since   
	the installation of the charger, at the moment charging stops.
	SessionStopEnergy-SessionStartEnergy is the amount of energy in Wh
	consumed in this charging session. \\
	28     & 2      & unknown4           & int16  & Always 0?                            \\
	30     & 2      & crc16              & uint16 & CRC-16/MODBUS over previous 30 bytes \\
\end{longtable}

\subsection{CCDT characteristics}
\label{ccdt_characteristics}

Table \ref{table.38} shwos wich characteristics of the \emph{CCDT} service are used.

\textbf{Note}: If the charger is an \emph{Nexxtender Home}, 
these characteristics are send under the \textbf{Generic} (C1) service. 
If the charger is an \emph{Nexxtender Mobile}, 
these characteristics are send under the \textbf{CCDT}(C5) service.

\begin{longtable}[]{@{}llllllll@{}}
	\caption{GATT CCDT service characteristics used}
	\label{table.38}
	\tabularnewline
	\toprule\noalign{}
	Handle & UUID & Name & R & W & WWR & N & Coding \\
	\midrule\noalign{}
	\endfirsthead
	\toprule\noalign{}
	Handle & UUID & Name & R & W & WWR & N & Coding \\
	\midrule\noalign{}
	\endhead
	\bottomrule\noalign{}
	\endlastfoot
	0x0022 & fd47416a-95fb-4206-88b5-b4a8045f75\textbf{c6} & \textbf{CCDT Command} &
	& P & & & HEX \\
	0x0024 & fd47416a-95fb-4206-88b5-b4a8045f75\textbf{c7} & \textbf{CCDT
	Status} & P & & & N & HEX \\
	0x0027 & fd47416a-95fb-4206-88b5-b4a8045f75\textbf{c8} & \textbf{CCDT
	Record} & P & & & & HEX \\
\end{longtable}

\textbf{CCDT Command}, \textbf{CCDT Status} and \textbf{CCDT Record}
work together to load transaction records from the device. 
The scenario is very similar to the one for \secref{cdr_characteristics}. 
The scenario is as shown in figure \ref{illustration.19}.

\begin{figure}[H]
	\caption{CCDT scenario}
	\label{illustration.19}
	\centering
	\includegraphics[scale=0.5]{CCDT.pdf}
\end{figure}

\begin{enumerate}
	
	\item
	      Client enables \textbf{CCDT Status} notifications.
	\item
	      Client reads the \textbf{CCDT Status} obtaining the new status,
	      indicating \emph{nnnnnnnn}, the number of CCDT records available. 
	      If \emph{nnnnnnnn} == 0, the scenario continues at point 6.
	\item
	      Client reads the \textbf{CCDT Record} obtaining the data of the record
	\item
	      Client writes the \textbf{CCDT Command}{[}0x01{]} for the next record.
	      Note that probably this also indicates that the record is successfully
	      stored by the Client, 
	      indicating that the Server can delete the record. 
	      So if you still need to sync with Nexxtmove.me, 
	      you should not perform above scenario with another app.
	\item
	      Client is notified of a \textbf{CCDT Status} change, 
	      containing the new status, indicating \emph{nnnnnnnn-\/-}, 
	      the new number of CDR records available. 
	      If \emph{nnnnnnnn} \textgreater{} 0, the scenario continues at point 3. 
	      Otherwise it is complete.
	\item
	      Client disables \textbf{CCDT Status} notifications.
\end{enumerate}

The number of CCDT records that can be stored by the Nexxtender Charger is limited. 
The exact limit is not known. 
If the limit is reached, it is expected that older records are overwritten.

\subsubsection{CCDT Command operation ids}
\label{ccdt_command_operation_ids}

Table \ref{table.57} lists the \emph{CCDT Command} operation
id\textquotesingle s.

\begin{longtable}[]{@{}lp{14cm}@{}}
	\caption{CCDT Command operation ids}
	\label{table.57}
	\tabularnewline
	\toprule\noalign{}
	Operation Id & Description                                                    \\
	\midrule\noalign{}
	\endfirsthead
	\toprule\noalign{}
	Operation Id & Description                                                    \\
	\midrule\noalign{}
	\endhead
	\bottomrule\noalign{}
	\endlastfoot
	0x01         & Read next record? Probably also means that the previous record 
	read from \textbf{CCDT Record} can be deleted from the Server. \\
	others?      &                                                                \\
\end{longtable}

\subsubsection{CCDT Status status}
\label{ccdt_status_status}

Possible \emph{CCDT Status status} values are shown in table \ref{table.58}.

\begin{longtable}[]{@{}ll@{}}
	\caption{CCDT Status status codes}
	\label{table.58}
	\tabularnewline
	\toprule\noalign{}
	Status Code       & Description                                     \\
	\midrule\noalign{}
	\endfirsthead
	\toprule\noalign{}
	Status Code       & Description                                     \\
	\midrule\noalign{}
	\endhead
	\bottomrule\noalign{}
	\endlastfoot
	0x\emph{nnnnnnnn} & \emph{nnnnnnnn} is an uint32 with the number of 
	records still to read \\
\end{longtable}

\subsubsection{CCDT Record operation (0xC4)}
\label{ccdt_record_operation_0xc4}

Returns a 16 bytes CCDT record related to a charging session that was summarized in a CDR record. 
Each CCDT record represents a status at some time within the start/stop timeframe of a CDR record. 
The 16 bytes are structured as shown in tables \ref{table.59} and \ref{table.60}

\begin{longtable}[]{@{}|l|l|l|l|l|l|l|l|l|l|l|l|l|l|l|l|@{}}
	\caption{CCDT Record data format overview}
	\label{table.59}\\  %<---- Problem
	\hline
	0 & 1 & 2 & 3 & 4 & 5 & 6 & 7 & 8 & 9 & 10 & 11 & 12 & 13 & 14 & 15 \\
	\hline
	\multicolumn{4}{@{}|l}{%
	timeStamp} & \multicolumn{4}{|l}{%
	eventEnergy} & \multicolumn{2}{|l|}{%
	quarterEnergy} & ccdtType & L1 & L2 & L3 & \multicolumn{2}{l|@{}}{%
	crc16} \\
	\hline
\end{longtable}

\hfill\break

\begin{longtable}[]{@{}llllp{10cm}@{}}
	\caption{CCDT Record data format detail}
	\label{table.60}
	\tabularnewline
	\toprule\noalign{}
	Offset & Length & Name          & Type   & Description                              \\
	\midrule\noalign{}
	\endfirsthead
	\toprule\noalign{}
	Offset & Length & Name          & Type   & Description                              \\
	\midrule\noalign{}
	\endhead
	\bottomrule\noalign{}
	\endlastfoot
	0      & 4      & timestamp     & uint32 & Timestamp of the record.                 \\
	4      & 4      & eventEnergy   & uint32 & Total amount of energy in Wh since the   
	installation of the charger, at the moment of the event. \\
	8      & 2      & quarterEnergy & uint16 & Total amount of energy in Wh since the   
	start of the current quarter. \\
	10     & 1      & ccdtType      & byte   & The type of the event according to table 
	\secref{table.64}.
	quarterEnergy, L1, L2 and L3 are 0 and eventEnergy does no longer
	increase for the events 0x21, 0x42, 0x60. \\
	11     & 1      & L1            & byte   & Car phase L1 current in A                \\
	12     & 1      & L2            & byte   & Car phase L2 current in A                \\
	13     & 1      & L3            & byte   & Car phase L3 current in A                \\
	14     & 2      & crc16         & uint16 & CRC-16/MODBUS over previous 13 bytes     \\
\end{longtable}

\hfill\break

\begin{longtable}[]{@{}lllp{11cm}@{}}
	\caption{CCDT type coding}
	\label{table.64}
	\tabularnewline
	\toprule\noalign{}
	Value & ASCII & Name          & Description                                      \\
	\midrule\noalign{}
	\endfirsthead
	\toprule\noalign{}
	Value & ASCII & Name          & Description                                      \\
	\midrule\noalign{}
	\endhead
	\bottomrule\noalign{}
	\endlastfoot
	0x21  & !     & CDR\_STARTED  & Seems to be the first CCDT record generated      
	within a CDR start/stop timeframe. \\
	0x41  & A     & NOT\_CHARGING & No longer charging. This occurs when the         
	charger is still charging, but the car is full and stopped the charging.
	A CCDT 0x41 event is generated every multiple of a quarter (hh:00,
	hh:15, hh:30, hh:45). \\
	0x42  & B     & CHARGE\_START & Charging has started                             \\
	0x49  & I     & CHARGE\_STOP  & Charging has stopped                             \\
	0x4A  & J     & CHARGING      & Still charging. During charging, a CCDT CHARGING 
	event is generated every multiple of a quarter (hh:00, hh:15, hh:30,
	hh:45). \\
	0x60  & `     & CDR\_STOPPED  & Seems to be the last CCDT record generated       
	within a CDR start/stop timeframe. \\
	x     & x     & Unknown       &                                                  \\
\end{longtable}

\subsection{Firmware characteristics}
\label{firmware_characteristics}

Table \ref{table.39} shows the characteristics of the \emph{Firmware} service that are used.

\textbf{Note}: If the charger is an \emph{Nexxtender Home}, 
these characteristics are send under the Generic (C1) service. 
If the charger is an \emph{Nexxtender Mobile}, 
these characteristics are send under the Firmware(C9) service.

\begin{longtable}[]{@{}llp{4cm}lllll@{}}
	\caption{GATT Firmware service characteristics used}
	\label{table.39}\\  %<---- Problem
	\tabularnewline
	\toprule\noalign{}
	Handle & UUID & Name & R & W & WWR & N & Coding \\
	\midrule\noalign{}
	\endfirsthead
	\toprule\noalign{}
	Handle & UUID & Name & R & W & WWR & N & Coding \\
	\midrule\noalign{}
	\endhead
	\bottomrule\noalign{}
	\endlastfoot
	0x0029 & fd47416a-95fb-4206-88b5-b4a8045f75\textbf{ca} &
	\textbf{Firmware Command} & & P? & & & \\
	0x002b & fd47416a-95fb-4206-88b5-b4a8045f75\textbf{cb} &
	\textbf{Firmware Status} & P & & & N & HEX \\
	0x002e & fd47416a-95fb-4206-88b5-b4a8045f75\textbf{cc} &
	\textbf{Firmware Wanted Chunk} & P & & & N & HEX \\
	0x0031 & fd47416a-95fb-4206-88b5-b4a8045f75\textbf{cd} &
	\textbf{Firmware Data Chunk} & & P? & & & HEX \\
\end{longtable}

\textbf{Firmware Command}, \textbf{Firmware Status}, \textbf{Firmware Wanted Chunk} 
and \textbf{Firmware Data Chunk} work together to load a new firmware on the device. 
The scenario is as shown in figure \ref{illustration.30}.

\begin{figure}[H]
	\caption{CCDT scenario}
	\label{illustration.30}
	\centering
	\includegraphics[scale=0.5]{FirmwareUpdate.pdf}
\end{figure}

\begin{enumerate}
	
	\item
	      Client enables \textbf{Firmware Status} notifications
	\item
	      Client reads the \textbf{Firmware Status} and extracts \emph{status}
	      and \emph{errorCode}.
	\item
	      As long as \emph{status} is not Firmware:IDLE(0x10) and \emph{errorCode} is not 0, 
	     the Client writes a Firmware:STOP(0x03) command to \textbf{Firmware Command} and
	\item
	      waits for a \textbf{Firmware Status} notification with a new
	      \emph{status} and \emph{errorCode}
	\item
	      As soon as \emph{status} is IDLE and \emph{errorCode} is 0, 
	      the Client writes a Firmware:START(0x02) command to \textbf{Firmware Command}
	      including the \emph{dataLength} of the new firmware.
	\item
	      The Client waits for a \textbf{Firmware Status} notification with a
	      new \emph{status} and \emph{errorCode}
	\item
	      If \emph{status} is Firmware:NEXT(0x12) and \emph{errorCode} is 0, 
	      the Client reads \textbf{Firmware Wanted Chunk} and extracts \emph{chunkNumber}, 
	      and
	\item
	      writes the data of the requested chunk to \textbf{Firmware Data Chunk}
	      and continues with step 6
	\item
	      If \emph{status} is Firmware:FINISHED(0x14) and \emph{errorCode} is 0,
	      the Client writes Firmware:RESET(0x04) to \textbf{Firmware Command}
	      and breaks out of the loop
	\item
	      If \emph{status} is not Firmware:NEXT(0x12) or Firmware:FINISHED(0x14)
	      or \emph{errorCode} is not 0, 
	      the client considers it as an error.
	\item
	      Client disables \textbf{Firmware Status} notifications
\end{enumerate}

\subsubsection{Firmware Command}
\label{firmware_command}

Writes an 8 byte command, structured as shown in tables \ref{table.69} and \ref{table.70}.

\begin{longtable}[]{@{}|l|l|l|l|l|l|l|l|@{}}
	\caption{Firmware Command data format overview}
	\label{table.69}\\  %<---- Problem
	\hline
	0 & 1 & 2 & 3 & 4 & 5 & 6 & 7 \\
	\hline
	operationId & \multicolumn{4}{|l}{%
	dataLength} & \multicolumn{3}{|l|}{%
	Unknown/RFU} \\
	\hline
\end{longtable}

\hfill\break

\begin{longtable}[]{@{}llllp{10cm}@{}}
	\caption{Firmware Command data format detail}
	\label{table.70}
	\tabularnewline
	\toprule\noalign{}
	Offset & Length & Name        & Type   & Description                              \\
	\midrule\noalign{}
	\endfirsthead
	\toprule\noalign{}
	Offset & Length & Name        & Type   & Description                              \\
	\midrule\noalign{}
	\endhead
	\bottomrule\noalign{}
	\endlastfoot
	0      & 1      & operationId & uint8  & Operation id according to table          
	\secref{table.65}. \\
	1      & 4      & dataLength  & uint32 & The total length of the firmware data in 
	bytes \\
	5      & 3      & Unknown/RFU &        &                                          \\
\end{longtable}

Table \ref{table.65} lists the \emph{Firmware Command} operation id\textquotesingle s.

\begin{longtable}[]{@{}llp{13cm}@{}}
	\caption{Firmware Command operation ids}
	\label{table.65}
	\tabularnewline
	\toprule\noalign{}
	Operation Id & Name  & Description                                            \\
	\midrule\noalign{}
	\endfirsthead
	\toprule\noalign{}
	Operation Id & Name  & Description                                            \\
	\midrule\noalign{}
	\endhead
	\bottomrule\noalign{}
	\endlastfoot
	0x02         & START & Starts a firmware loading sequence                     \\
	0x03         & STOP  & Stops a firmware loading sequence. Firmware loading is aborted? \\
	0x04         & RESET & Resets the charger after a successful firmware loading sequence. 
    The charger restarts with the new firmware. \\
	others?      &       &                                                        \\
\end{longtable}

\subsubsection{Firmware Status}
\label{firmware_status}

Returns an 8 byte status, structured as shown in tables \ref{table.67} and \ref{table.68}.

\begin{longtable}[]{@{}|l|l|l|l|l|l|l|l|@{}}
	\caption{Firmware Status data format overview}
	\label{table.67}\\  %<---- Problem
	\hline
	0 & 1 & 2 & 3 & 4 & 5 & 6 & 7 \\
	\hline
	status & errorCode & \multicolumn{6}{l|@{}}{%
	Unknown/RFU} \\
	\hline
\end{longtable}

\hfill\break

\begin{longtable}[]{@{}lllll@{}}
	\caption{Firmware Status data format detail}
	\label{table.68}
	\tabularnewline
	\toprule\noalign{}
	Offset & Length & Name        & Type  & Description                                  \\
	\midrule\noalign{}
	\endfirsthead
	\toprule\noalign{}
	Offset & Length & Name        & Type  & Description                                  \\
	\midrule\noalign{}
	\endhead
	\bottomrule\noalign{}
	\endlastfoot
	0      & 1      & status      & uint8 & Status of the firmware download according to 
	table \secref{table.66} \\
	1      & 1      & errorCode   & uint8 & Error code. 0x00 is success                  \\
	2      & 6      & Unknown/RFU &       &                                              \\
\end{longtable}

\hfill\break

Table \ref{table.66} shows the coding of the \emph{status} field.

\begin{longtable}[]{@{}llp{12cm}@{}}
	\caption{Status coding}
	\label{table.66}
	\tabularnewline
	\toprule\noalign{}
	Status Code & Name     & Description                                              \\
	\midrule\noalign{}
	\endfirsthead
	\toprule\noalign{}
	Status Code & Name     & Description                                              \\
	\midrule\noalign{}
	\endhead
	\bottomrule\noalign{}
	\endlastfoot
	0x00        & BUSY     & Charger is busy processing a \textbf{Firmware Command}?  \\
	0x10        & IDLE     & Charger is not busy with a firmware downloading session  \\
	0x11        & CRC      & ?                                                        \\
	0x12        & NEXT     & The charger is ready to receive the next chunk. 
	Read the wanted chunk number from \textbf{Firmware Wanted Chunk} and 
	write the 	corresponding chunk data to \textbf{Firmware Data Chunk} \\
	0x14        & FINISHED & The complete firmware is loaded to the charger. 
	If the   errorCode is 0x00, restart the charger with the new firmware using the
	\textbf{Firmware Command}{[}Firmware:RESET{]}. \\
	Others      & UNKNOWN  &                                                          \\
\end{longtable}

\subsubsection{Firmware Wanted Chunk}
\label{firmware_wanted_chunk}

Returns a 2 byte status, structured as shown in tables \ref{table.71} and \ref{table.72}.

\begin{longtable}[]{|l|l|}
	\caption{Firmware Wanted Chunk data format overview}
	\label{table.71}\\  %<---- Problem
	\hline
	0 & 1 \\
	\hline
	\multicolumn{2}{|l|}{%
	chunkNumber} \\
	\hline
\end{longtable}

\hfill\break

\begin{longtable}[]{@{}lllll@{}}
	\caption{Firmware Wanted Chunk data format detail}
	\label{table.72}
	\tabularnewline
	\toprule\noalign{}
	Offset & Length & Name        & Type   & Description                       \\
	\midrule\noalign{}
	\endfirsthead
	\toprule\noalign{}
	Offset & Length & Name        & Type   & Description                       \\
	\midrule\noalign{}
	\endhead
	\bottomrule\noalign{}
	\endlastfoot
	0      & 2      & chunkNumber & uint16 & Number of the firmware data chunk 
	expected \\
\end{longtable}

\subsubsection{Firmware Data Chunk}
\label{firmware_data_chunk}

Sends a 16-byte or 64-byte chunk of the firmware. 
If the charger is an \emph{Nexxtender Home}, the chunk size is 64. 
If the charger is an \emph{Nexxtender Mobile}, the chunk size is 16.
The bytes are structured as shown in table \ref{table.73}.

\begin{longtable}[]{@{}llllp{10cm}@{}}
	\caption{Firmware Data Chunk data format detail}
	\label{table.73}
	\tabularnewline
	\toprule\noalign{}
	Offset & Length   & Name  & Type                & Description             \\
	\midrule\noalign{}
	\endfirsthead
	\toprule\noalign{}
	Offset & Length   & Name  & Type                & Description             \\
	\midrule\noalign{}
	\endhead
	\bottomrule\noalign{}
	\endlastfoot
	0      & 16 or 64 & chunk & uint8{[}16 or 64{]} & The firmware data chunk corresponding 
	with the chunkNumber returned in \textbf{Firmware Wanted Chunk} \\
\end{longtable}

\subsection{Charging characteristics}
\label{charging_characteristics}

Table \ref{table.40} shows the characteristics of the \emph{Charging} service that are used.

\textbf{Note}: If the charger is an \emph{Nexxtender Home}, 
these characteristics are send under the \textbf{Generic} (C1) service. 
If the charger is an \emph{Nexxtender Mobile}, these characteristics are send
under the \textbf{Charging} (CE) service.

\begin{longtable}[]{@{}llp{3.7cm}lllll@{}}
	\caption{Charging service characteristics used}
	\label{table.40}
	\tabularnewline
	\toprule\noalign{}
	Handle & UUID & Name & R & W & WWR & N & Coding \\
	\midrule\noalign{}
	\endfirsthead
	\toprule\noalign{}
	Handle & UUID & Name & R & W & WWR & N & Coding \\
	\midrule\noalign{}
	\endhead
	\bottomrule\noalign{}
	\endlastfoot
	0x0033 & fd47416a-95fb-4206-88b5-b4a8045f75\textbf{cf} &
	\textbf{Charging Basic Data} & P & & & N & HEX \\
	0x0036 & fd47416a-95fb-4206-88b5-b4a8045f75\textbf{d0} &
	\textbf{Charging Grid Data} & P & & & N & HEX \\
	0x0039 & fd47416a-95fb-4206-88b5-b4a8045f75\textbf{da} &
	\textbf{Charging Car Data} & P & & & N & HEX \\
	0x003f & fd47416a-95fb-4206-88b5-b4a8045f75\textbf{db} &
	\textbf{Charging Advanced Data} & P & & & N & HEX \\
\end{longtable}

The \emph{Nexxtender Home} supports all 4 characteristics. 
The \emph{Nexxtender Mobile} only supports the \textbf{Charging Basic Data} characteristic.

\subsubsection{Charging Basic Data (0xCF)}
\label{charging_basic_data_0xcf}

Returns 14 bytes, structured as shown in tables \ref{table.41} and \ref{table.42}.

\begin{longtable}[]{@{}|l|l|l|l|l|l|l|l|l|l|l|l|l|l|@{}}
	\caption{Charging Basic Data data format overview}
	\label{table.41}\\  %<---- Problem
	\hline
	0 & 1 & 2 & 3 & 4 & 5 & 6 & 7 & 8 & 9 & 10 & 11 & 12 & 13 \\
	\hline
	\multicolumn{2}{@{}|l|}{%
	seconds} & discriminator & status & \multicolumn{4}{l}{%
	RFU} & \multicolumn{4}{|l|}{%
	energy} & RFU & phaseCount \\
	\hline
\end{longtable}

\hfill\break

\begin{longtable}[]{@{}llllp{10cm}@{}}
	\caption{Charging Basic Data data format detail}
	\label{table.42}
	\tabularnewline
	\toprule\noalign{}
	Offset & Length & Name          & Type   & Description                            \\
	\midrule\noalign{}
	\endfirsthead
	\toprule\noalign{}
	Offset & Length & Name          & Type   & Description                            \\
	\midrule\noalign{}
	\endhead
	\bottomrule\noalign{}
	\endlastfoot
	0      & 2      & seconds       & int16  & number of seconds since start of       
	charging\footnote{What is exact definition?} \\
	2      & 1      & discriminator & uint8  & Discriminator state according to table \secref{table.43} \\
	3      & 1      & status        & uint8  & Status according to table	\secref{table.44} \\
	4      & 4      & RFU           & uint32 & RFU                                    \\
	8      & 4      & energy        & uint32 & Total energy in Wh charged during this session? \\
	12     & 1      & RFU           & uint8  & RFU                                    \\
	13     & 1      & phaseCount    & uint8  & Charging Phase Count                   \\
\end{longtable}

\hfill\break

Table \ref{table.43} shows the coding for the \emph{discriminator} field.

\begin{longtable}[]{@{}lllllllllll@{}}
	\caption{Discriminator coding}
	\label{table.43}
	\tabularnewline
	\toprule\noalign{}
	b7 & b6 & b5 & b4 & b3 & b2 & b1 & b0 & Value & Name     & Description      \\
	\midrule\noalign{}
	\endfirsthead
	\toprule\noalign{}
	b7 & b6 & b5 & b4 & b3 & b2 & b1 & b0 & Value & Name     & Description      \\
	\midrule\noalign{}
	\endhead
	\bottomrule\noalign{}
	\endlastfoot
	0  & 0  & 0  & 0  & 0  & 0  & 0  & 1  & 1     & STARTED  &                  \\
	0  & 0  & 0  & 0  & 0  & 0  & 1  & 0  & 2     & CHARGING &                  \\
	0  & 0  & 0  & 0  & 0  & 0  & 1  & 1  & 3     & STOPPED  &                  \\
	x  & x  & x  & x  & x  & x  & x  & x  & x     & UNKNOWN  & Any other value: 
	UNKNOWN \\
\end{longtable}

\hfill\break

Table \ref{table.44} shows the coding for the \emph{status} field.

\begin{longtable}[]{@{}llll@{}}
	\caption{Status coding}
	\label{table.44}
	\tabularnewline
	\toprule\noalign{}
	Value & ASCII & Name     & Description              \\
	\midrule\noalign{}
	\endfirsthead
	\toprule\noalign{}
	Value & ASCII & Name     & Description              \\
	\midrule\noalign{}
	\endhead
	\bottomrule\noalign{}
	\endlastfoot
	0x42  & B     & PLUGGED  &                          \\
	0x43  & C     & CHARGING &                          \\
	0x44  & D     & CHARGING &                          \\
	0x45  & E     & FAULT    &                          \\
	0x46  & F     & FAULT    &                          \\
	x     & x     & UNKNOWN  & Any other value: UNKNOWN \\
\end{longtable}
\hfill\break

If \textbf{Charging Basic Data} is set up to notify the client of changes, 
the scenario is as shown in figure \ref{illustration.12}.

\begin{figure}[H]
	\caption{Charging Basic Data notification}
	\label{illustration.12}
	\centering
	\includegraphics[scale=0.5]{ChargingBasicDataNotify.pdf}
\end{figure}

Experience shows that there is a notification approximately every 5 seconds.

\subsubsection{Charging Grid Data (0xD0)}
\label{charging_grid_data_0xd0}

Returns 16 bytes related to the current electrical current on the grid.
This is the current as measured by the 3 Current Transformers (CT). 
The 16 bytes are structured as showns in tables \ref{table.45} and \ref{table.46}.

\begin{longtable}[]{@{}|l|l|l|l|l|l|l|l|l|l|l|l|l|l|l|l|@{}}
	\caption{Charging Grid Data data format overview}
	\label{table.45}\\  %<---- Problem
	\hline
	0 & 1 & 2 & 3 & 4 & 5 & 6 & 7 & 8 & 9 & 10 & 11 & 12 & 13 & 14 & 15 \\
	\hline
	\multicolumn{4}{@{}|l}{%
	timestamp} & \multicolumn{2}{|l}{%
	L1} & \multicolumn{2}{|l}{%
	L2} & \multicolumn{2}{|l}{%
	L3} & \multicolumn{2}{|l}{%
	consumed} & \multicolumn{2}{|l}{%
	interval} & \multicolumn{2}{|l|@{}}{%
	crc16} \\
	\hline
\end{longtable}
\hfill\break

\begin{longtable}[]{@{}llllp{10cm}@{}}
	\caption{Charging Grid Data data format detail}
	\label{table.46}
	\tabularnewline
	\toprule\noalign{}
	Offset & Length & Name      & Type   & Description                                                                               \\
	\midrule\noalign{}
	\endfirsthead
	\toprule\noalign{}
	Offset & Length & Name      & Type   & Description                                                                               \\
	\midrule\noalign{}
	\endhead
	\bottomrule\noalign{}
	\endlastfoot
	0      & 4      & timestamp & uint32 & Current time in UTC Unix Time                                                             \\
	4      & 2      & L1        & int16  & Grid phase L1 current in                                                                  
	dA\footnote{\label{fn:negative}A negative value indicates that a power surplus is delivered to the Grid. 
	That can only happen in case there are solar panels and
	solar power is more than what is consumed by the house and the EV} \\
	6      & 2      & L2        & int16  & Grid phase L2 current in dA\footref{fn:negative} \\
	8      & 2      & L3        & int16  & Grid phase L3 current in dA\footref{fn:negative} \\
	10     & 2      & consumed  & int16  & Total power injected in the Grid in cWh during this quarter? See remark after this table. \\
	12     & 2      & interval  & uint16 & Counts the seconds in the current quarter.                                                
	Is set to 0 every multiple of a quarter (hh:00, hh:15, hh:30, hh:45). \\
	14     & 2      & crc16     & uint16 & CRC-16/MODBUS over previous 14 bytes                                                      \\
\end{longtable}
\hfill\break

It is unclear what the \emph{consumed} field is.
From the name one would expect it is the amount of energy consumed in Wh from the grid during the current quarter. 
But that does not seem to be true, at least not for firmware 2.53.2.
The following is observed while inspecting Grid Data in 5 second intervals.

\begin{itemize}
	\item \emph{consumed}  is reset to 0 at he start of each quarter, i.e. when \emph{interval} is also set to 0.
	\item \emph{consumed} never decreases.
	\item \emph{consumed} does not change as long as L1 remains positive, i.e. net grid consumption. 
	This is the opposed of what the name suggests. 
	That also means that for installations without solar panels, \emph{consumed} always stays 0. 
	\item \emph{consumed} increases if power is delivered to the grid. 
	With a car being charged in ECO mode, so consuming all power generated by solar panels, the increase is slow. 
	If the car is not charging while the solar panels produce power, the increase is fast.
	\item \emph{consumed} seems to express the power in cWh.
	\item Although \emph{consumed} is coded on 16 bits, only 12 seem to be used. 
	The value increases to about 3276, after which it jumps to around -3276 and keeps on increasing (so towards positive values). 
	An example is a case where after incrementing each 5 seconds with about 0x80 until the value 0x0C8F,
	the next value is 0xF37B, which is a negative value. 
	This usually only happens when the car is not charging and there is a lot of solar energy available. 
	Otherwise, most of the solar energy is consumed by the charger and \emph{consumed} only increments a little; 
	insufficient to cause the overflow.
\end{itemize}

If \textbf{Charging Grid Data} is set up to notify the client of changes, 
the scenario is as shown in figure \ref{illustration.13}.

\begin{figure}[H]
	\caption{Charging Grid Data notification}
	\label{illustration.13}
	\centering
	\includegraphics[scale=0.5]{ChargingGridDataNotify.pdf}
\end{figure}

Experience shows that there is a notification approximately every 5 seconds.

\subsubsection{Charging Car Data (0xDA)}
\label{charging_car_data_0xda}

Returns 18 bytes related to the current electrical current towards the car. 
This is the current as provided by the charger to the car. 
The 18 bytes are structured as shown in tables \ref{table.47} and \ref{table.48}.

\begin{longtable}[]{@{}|l|l|l|l|l|l|l|l|l|l|l|l|l|l|l|l|l|l|@{}}
	\caption{Charging Car Data data format overview}
	\label{table.47}\\  %<---- Problem
	\hline
	0 & 1 & 2 & 3 & 4 & 5 & 6 & 7 & 8 & 9 & 10 & 11 & 12 & 13 & 14 & 15 & 16 
	& 17 \\
	\hline
	\multicolumn{4}{@{}|l}{%
	timestamp} & \multicolumn{2}{|l}{%
	L1} & \multicolumn{2}{|l}{%
	L2} & \multicolumn{2}{|l}{%
	L3} & \multicolumn{2}{|l}{%
	P1} & \multicolumn{2}{|l}{%
	P2} & \multicolumn{2}{|l}{%
	P3} & \multicolumn{2}{|l|@{}}{%
	crc16} \\
	\hline
\end{longtable}

\hfill\break

\begin{longtable}[]{@{}lllll@{}}
	\caption{Charging Car Data data format detail}
	\label{table.48}
	\tabularnewline
	\toprule\noalign{}
	Offset & Length & Name      & Type   & Description                            \\
	\midrule\noalign{}
	\endfirsthead
	\toprule\noalign{}
	Offset & Length & Name      & Type   & Description                            \\
	\midrule\noalign{}
	\endhead
	\bottomrule\noalign{}
	\endlastfoot
	0      & 4      & timestamp & uint32 & Current time in UTC Unix Time          \\
	4      & 2      & L1        & int16  & Car phase L1 current in dA             \\
	6      & 2      & L2        & int16  & Car phase L2 current in dA             \\
	8      & 2      & L3        & int16  & Car phase L3 current in dA             \\
	10     & 2      & P1        & int16  & Car phase L1 power consumption in Watt \\
	12     & 2      & P2        & int16  & Car phase L2 power consumption in Watt \\
	14     & 2      & P3        & int16  & Car phase L2 power consumption in Watt \\
	16     & 2      & crc16     & uint16 & CRC-16/MODBUS over previous 16 bytes   \\
\end{longtable}
\hfill\break
If \textbf{Charging Car Data} is set up to notify the client of changes,
the scenario is as shown in figure \ref{illustration.14}.

\begin{figure}[H]
	\caption{Charging Car Data notification}
	\label{illustration.14}
	\centering
	\includegraphics[scale=0.5]{ChargingCarDataNotify.pdf}
\end{figure}
\hfill\break

Experience shows that there is a notification approximately every 5 seconds.

\subsubsection{Charging Advanced Data (0xDB)}
\label{charging_advanced_data_0xdb}

Returns 18 bytes related to the advanced charging data.  
The 18 bytes are structured as shown in tables \ref{table.49} and \ref{table.50}.

\begin{longtable}[]{@{}|l|l|l|l|l|l|l|l|l|l|l|l|l|l|l|l|l|l|@{}}
	\caption{Charging Advanced Data data format overview}
	\label{table.49}\\  %<---- Problem
	\hline
	0 & 1 & 2 & 3 & 4 & 5 & 6 & 7 & 8 & 9 & 10 & 11 & 12 & 13 & 14 & 15 & 16 
	& 17 \\
	\hline
	\multicolumn{4}{@{}|l}{%
	timestamp} & \multicolumn{2}{|l}{%
	iAvailable} & \multicolumn{4}{|l}{%
	gridPower} & \multicolumn{4}{|l|}{%
	carPower} & authorizationStatus & errorCode & \multicolumn{2}{l|@{}}{%
	crc16} \\
	\hline
\end{longtable}
\hfill\break

\begin{longtable}[]{@{}llllp{9cm}@{}}
	\caption{Charging Advanced Data data format detail}
	\label{table.50}
	\tabularnewline
	\toprule\noalign{}
	Offset & Length & Name                & Type   & Description                                 \\
	\midrule\noalign{}
	\endfirsthead
	\toprule\noalign{}
	Offset & Length & Name                & Type   & Description                                 \\
	\midrule\noalign{}
	\endhead
	\bottomrule\noalign{}
	\endlastfoot
	0      & 4      & timestamp           & uint32 & Current time in UTC Unix Time               \\
	4      & 2      & iAvailable          & int16  & Available current in A                      \\
	6      & 4      & gridPower           & int32  & Total power consumption from the grid in Watt\footref{fn:negative}. 
	\footnote{What is the difference with consumed in \textbf{Charging Grid Data}} 
	Equals P\textsubscript{grid} in section \secref{ecosystem_overview}. \\
	10     & 4      & carPower            & int32  & Total power consumption from the car in Watt. 
	\footnote{Is this the sum of Car P1, P2 and P3 in \textbf{Charging Car Data}?} 
	Equals P\textsubscript{EV} in section \secref{ecosystem_overview}. \\
	14     & 1      & authorizationStatus & uint8  & Authorization status according to	table	\secref{table.52} \\
	15     & 1      & errorCode           & uint8  & Error code returned by the \emph{Nexxtender Home}. 
	Values unknown. \\
	16     & 2      & crc16               & uint16 & CRC-16/MODBUS over previous 16 bytes        \\
\end{longtable}

Tables \ref{table.51} and \ref{table.52} show the coding for the \emph{authorizationStatus bit} 
field and some special combinations.

\begin{longtable}[]{@{}lllllllllll@{}}
	\caption{authorizationStatus coding}
	\label{table.51}
	\tabularnewline
	\toprule\noalign{}
	b7 & b6 & b5 & b4 & b3 & b2 & b1 & b0 & Value & Name                  & Description \\
	\midrule\noalign{}
	\endfirsthead
	\toprule\noalign{}
	b7 & b6 & b5 & b4 & b3 & b2 & b1 & b0 & Value & Name                  & Description \\
	\midrule\noalign{}
	\endhead
	\bottomrule\noalign{}
	\endlastfoot
	x  & x  & x  & x  & x  & x  & x  & 1  & 0x01  & Unauthorized          &             \\
	x  & x  & x  & x  & x  & x  & 1  & x  & 0x02  & Authorized default    &             \\
	x  & x  & x  & x  & x  & 1  & x  & x  & 0x04  & Charge stopped in app &             \\
	x  & x  & x  & x  & 1  & x  & x  & x  & 0x08  & RFU1                  &             \\
	x  & x  & x  & 1  & x  & x  & x  & x  & 0x10  & Charge paused         &             \\
	x  & x  & 1  & x  & x  & x  & x  & x  & 0x20  & MAX                   &             \\
	x  & 1  & x  & x  & x  & x  & x  & x  & 0x40  & ECO                   &             \\
	1  & x  & x  & x  & x  & x  & x  & x  & 0x80  & RFU2                  &             \\
\end{longtable}
\hfill\break

\begin{longtable}[]{@{}lllllllllll@{}}
	\caption{authorizationStatus values}
	\label{table.52}
	\tabularnewline
	\toprule\noalign{}
	b7 & b6 & b5 & b4 & b3 & b2 & b1 & b0 & Value & Name               & Description \\
	\midrule\noalign{}
	\endfirsthead
	\toprule\noalign{}
	b7 & b6 & b5 & b4 & b3 & b2 & b1 & b0 & Value & Name               & Description \\
	\midrule\noalign{}
	\endhead
	\bottomrule\noalign{}
	\endlastfoot
	0  & 0  & 0  & 0  & 0  & 0  & 0  & 1  & 0x01  & Unauthorized       &             \\
	0  & 0  & 0  & 0  & 0  & 0  & 1  & 0  & 0x02  & Authorized default &             \\
	0  & 0  & 1  & 0  & 0  & 0  & 1  & 0  & 0x22  & Authorized MAX     &             \\
	0  & 1  & 0  & 0  & 0  & 0  & 1  & 0  & 0x42  & Authorized ECO     &             \\
\end{longtable}
\hfill\break

If \textbf{Charging Advanced Data} is set up to notify the client of changes, 
the scenario is as shown in figure \ref{illustration.15}.

\begin{figure}[H]
	\caption{Charging Advanced Data notification}
	\label{illustration.15}
	\centering
	\includegraphics[scale=0.5]{ChargingAdvancedDataNotify.pdf}
\end{figure}

Experience shows that there is a notification approximately every 5 seconds.

\nocite{*}          % include all references
\bibliographystyle{plain} % We choose the "plain" reference style
\bibliography{nexxtender.bib} % Entries are in the nexxtender.bib file


\end{document}